\documentclass[12pt]{report}
\usepackage{amsmath, amssymb, amsthm}
\usepackage{graphicx}
\usepackage{caption, subcaption}
\usepackage{geometry}
\usepackage{hyperref}
\usepackage{booktabs}
\usepackage{float}
\usepackage{natbib}
\usepackage{setspace}
\usepackage{tikz}
\usepackage{algorithm}
\usepackage{algorithmic}
\usepackage{mathrsfs}
\usepackage{bbm}

\geometry{
    a4paper,
    left=25mm,
    right=25mm,
    top=25mm,
    bottom=25mm,
}

% Theorem environments
\newtheorem{theorem}{Theorem}[section]
\newtheorem{lemma}[theorem]{Lemma}
\newtheorem{proposition}[theorem]{Proposition}
\newtheorem{corollary}[theorem]{Corollary}
\theoremstyle{definition}
\newtheorem{definition}[theorem]{Definition}
\newtheorem{remark}[theorem]{Remark}

% Custom commands
\newcommand{\E}{\mathbb{E}}
\newcommand{\R}{\mathbb{R}}
\newcommand{\N}{\mathcal{N}}
\newcommand{\D}{\mathcal{D}}
\newcommand{\KL}{\text{KL}}
\newcommand{\loss}{\mathcal{L}}
\newcommand{\norm}[1]{\left\| #1 \right\|}
\DeclareMathOperator{\Tr}{Tr}
\DeclareMathOperator{\diag}{diag}

\title{Compact Diffusion Models: A Comprehensive Study of Knowledge Distillation, Pruning, Quantization, and Symmetry-Aware Architectures}
\author{Ajit Kumar Sorout}
\date{\today}

\begin{document}

\begin{titlepage}
    \centering
    \vspace*{3cm}
    {\huge\bfseries Compact Diffusion Models:\par}
    \vspace{0.5cm}
    {\LARGE\bfseries A Comprehensive Study of Knowledge Distillation, Pruning, Quantization, and Symmetry-Aware Architectures\par}
    \vspace{2cm}
    {\Large\itshape Ajit Kumar \\ Pronay Dutta\par}
    \vfill
    {\large \today\par}
\end{titlepage}

\pagenumbering{roman}
\begin{abstract}
Diffusion models have emerged as state-of-the-art generative models, achieving unprecedented success in high-fidelity image synthesis by learning to reverse a stochastic noise corruption process. However, their deployment on resource-constrained platforms remains challenging due to substantial computational and memory requirements. This comprehensive report presents a rigorous mathematical framework for developing compact diffusion models through the strategic integration of knowledge distillation, structured pruning, quantization, and symmetry-aware architectural design.

We provide detailed theoretical foundations for each compression technique, including complete derivations of training objectives, convergence properties, and performance bounds. Our experimental implementation on CIFAR-10, while limited by computational constraints to 10,000 images and approximately 200 training epochs, demonstrates the practical feasibility of achieving significant model compression (74\% size reduction) with controlled performance degradation. Crucially, we introduce a novel theoretical framework for incorporating physical symmetries and invariances into diffusion model architectures, drawing from principles in geometric deep learning, equivariant neural networks, and physics-informed machine learning. This symmetry-aware perspective not only enhances data efficiency and model compactness but also ensures that generated samples respect fundamental conservation laws and geometric constraints relevant to specific application domains.

The report concludes with a thorough analysis of current limitations, including the impact of aggressive compression on Fréchet Inception Distance (FID) scores, and outlines concrete directions for future research in latent diffusion models, advanced quantization schemes, and symmetry-preserving sampling algorithms.
\end{abstract}

\tableofcontents
\listoffigures
\listoftables

\newpage
\pagenumbering{arabic}

\chapter{Introduction}

\section{Motivation and Context}

The field of generative modeling has undergone a paradigm shift with the advent of diffusion probabilistic models~\citep{sohl2015deep, ho2020denoising, song2020score}. These models learn to generate high-dimensional data by iteratively denoising samples from a simple prior distribution, typically an isotropic Gaussian. Unlike generative adversarial networks (GANs)~\citep{goodfellow2014generative} or variational autoencoders (VAEs)~\citep{kingma2013auto}, diffusion models offer stable training dynamics, mode coverage guarantees, and interpretable generation processes grounded in stochastic differential equations (SDEs) and thermodynamic principles.

Despite these advantages, diffusion models present significant computational challenges:

\begin{enumerate}
    \item \textbf{Iterative Generation}: Sampling requires sequential evaluation through hundreds or thousands of diffusion timesteps, leading to prohibitive inference latency.
    \item \textbf{Model Scale}: State-of-the-art diffusion models often contain hundreds of millions of parameters, demanding substantial memory resources.
    \item \textbf{Training Complexity}: Learning the score function across the entire diffusion trajectory necessitates extensive computational resources and training time.
\end{enumerate}

These constraints severely limit deployment scenarios, particularly for edge devices, mobile platforms, and real-time applications. Addressing this challenge requires a multifaceted approach combining model compression, architectural innovations, and algorithm-level optimizations.

\section{Compression Paradigms for Neural Networks}

Model compression techniques have been extensively studied in the context of discriminative models~\citep{hinton2015distilling, han2015deep, jacob2018quantization}. These methods broadly fall into four categories:

\begin{description}
    \item[Knowledge Distillation] Transferring learned representations from a large teacher network to a compact student network by minimizing divergence between their output distributions~\citep{hinton2015distilling}.
    \item[Pruning] Systematically removing redundant or less important parameters based on magnitude, gradient, or higher-order sensitivity metrics~\citep{han2015learning, molchanov2016pruning}.
    \item[Quantization] Reducing numerical precision of weights and activations from 32-bit floating-point to lower bit-widths (e.g., 8-bit integers)~\citep{jacob2018quantization, krishnamoorthi2018quantizing}.
    \item[Architectural Search] Automatically discovering efficient network topologies optimized for specific compute/accuracy trade-offs~\citep{tan2019efficientnet}.
\end{description}

While these techniques have proven effective for convolutional and transformer-based discriminative models, their application to generative diffusion models introduces unique challenges related to maintaining sample quality, preserving mode coverage, and ensuring stable multi-step generation processes.

\section{Symmetry and Inductive Bias}

A parallel development in machine learning has been the recognition that incorporating domain-specific symmetries and invariances as inductive biases can dramatically improve sample efficiency, generalization, and model compactness~\citep{cohen2016group, bronstein2021geometric}. In physics-informed machine learning, respecting fundamental symmetries (e.g., Galilean invariance, rotational equivariance, gauge symmetry) ensures that learned representations align with physical laws~\citep{cranmer2020discovering}.

For diffusion models, symmetry-aware architectures offer several advantages:

\begin{itemize}
    \item \textbf{Parameter Efficiency}: Weight sharing induced by group equivariance reduces effective parameter count.
    \item \textbf{Data Efficiency}: Models that respect symmetries require fewer training examples to generalize.
    \item \textbf{Physical Consistency}: Generated samples automatically satisfy conservation laws and geometric constraints.
    \item \textbf{Compositional Generation}: Symmetries enable factorization of generation process across independent modes.
\end{itemize}

\section{Objectives and Contributions}

The primary objectives of this work are:

\begin{enumerate}
    \item \textbf{Theoretical Framework}: Establish rigorous mathematical foundations for each compression technique in the context of diffusion models, including detailed derivations of training objectives and performance bounds.
    
    \item \textbf{Empirical Validation}: Demonstrate practical feasibility through implementation and evaluation on CIFAR-10, analyzing trade-offs between compression ratio, inference speed, and generation quality.
    
    \item \textbf{Symmetry Integration}: Develop a comprehensive framework for incorporating physical symmetries into diffusion architectures, with concrete proposals for equivariant denoising networks.
    
    \item \textbf{Analysis and Future Directions}: Critically evaluate current limitations and propose concrete research directions for advancing compact diffusion models.
\end{enumerate}

Our contributions include:

\begin{itemize}
    \item Complete mathematical derivations for diffusion model training objectives, including variational lower bounds and score matching formulations.
    \item Detailed analysis of knowledge distillation for diffusion models, with novel loss formulations incorporating intermediate feature matching and temporal consistency.
    \item Rigorous treatment of structured pruning with dependency graph analysis and layer-wise sensitivity metrics.
    \item Comprehensive coverage of quantization-aware training with gradient approximation techniques.
    \item Novel theoretical framework for symmetry-aware diffusion models, connecting concepts from geometric deep learning to generative modeling.
    \item Experimental results demonstrating 74\% model size reduction with controlled performance degradation.
\end{itemize}

\section{Document Organization}

This report is organized as follows:

\begin{description}
    \item[Chapter 2] provides comprehensive mathematical background on diffusion probabilistic models, including forward/reverse processes, training objectives, and connections to score-based generative modeling and SDEs.
    
    \item[Chapter 3] presents detailed theoretical foundations for model compression techniques: knowledge distillation, pruning, and quantization, with complete derivations and convergence analyses.
    
    \item[Chapter 4] discusses practical training techniques including mixed precision training, exponential moving averages, and gradient scaling strategies.
    
    \item[Chapter 5] describes the conditional U-Net architecture employed in our implementation, with emphasis on design choices for compact models.
    
    \item[Chapter 6] introduces a comprehensive framework for symmetry-aware diffusion models, covering group theory fundamentals, equivariant architectures, and physics-informed constraints.
    
    \item[Chapter 7] presents experimental methodology, results, and critical analysis of our compact diffusion model implementation.
    
    \item[Chapter 8] concludes with discussion of limitations and future research directions.
\end{description}

\chapter{Mathematical Foundations of Diffusion Models}

This chapter establishes the rigorous mathematical framework underlying diffusion probabilistic models. We begin with the fundamental forward and reverse diffusion processes, derive the variational lower bound training objective, and explore connections to score-based generative modeling and stochastic differential equations.

\section{Probabilistic Framework}

\subsection{Forward Diffusion Process}

Let $\mathbf{x}_0 \sim q_{\text{data}}(\mathbf{x}_0)$ denote a sample from the true data distribution. The forward diffusion process is a discrete-time Markov chain that gradually corrupts $\mathbf{x}_0$ by adding Gaussian noise over $T$ timesteps:

\begin{equation}
q(\mathbf{x}_t | \mathbf{x}_{t-1}) = \N(\mathbf{x}_t; \sqrt{1 - \beta_t} \mathbf{x}_{t-1}, \beta_t \mathbf{I}),
\label{eq:forward_step}
\end{equation}

where $\{\beta_t\}_{t=1}^T$ is a variance schedule satisfying $0 < \beta_t < 1$ for all $t$.

\begin{figure}[ht]
    \centering
    \includegraphics[width=0.9\linewidth]{diffusion_forward.png}
    \caption{Illustration of the forward diffusion process progressively adding Gaussian noise to data samples over $T$ timesteps.}
    \label{fig:diffusion_forward}
\end{figure}

\subsubsection{Reparameterization and Closed Form}

Using the reparameterization trick, we can express Equation~\eqref{eq:forward_step} as:

\begin{equation}
\mathbf{x}_t = \sqrt{1 - \beta_t} \mathbf{x}_{t-1} + \sqrt{\beta_t} \boldsymbol{\epsilon}_{t-1}, \quad \boldsymbol{\epsilon}_{t-1} \sim \N(\mathbf{0}, \mathbf{I}).
\label{eq:forward_reparam}
\end{equation}

A key property of Gaussian diffusion processes is that the marginal distribution $q(\mathbf{x}_t | \mathbf{x}_0)$ can be computed in closed form. Define:

\begin{align}
\alpha_t &= 1 - \beta_t, \label{eq:alpha_def} \\
\bar{\alpha}_t &= \prod_{s=1}^t \alpha_s. \label{eq:alpha_bar_def}
\end{align}

Then by the Markov property and Gaussian closure:

\begin{proposition}[Closed-Form Forward Process]
For any timestep $t \in \{1, \ldots, T\}$:
\begin{equation}
q(\mathbf{x}_t | \mathbf{x}_0) = \N(\mathbf{x}_t; \sqrt{\bar{\alpha}_t} \mathbf{x}_0, (1 - \bar{\alpha}_t) \mathbf{I}).
\label{eq:forward_closed}
\end{equation}
\end{proposition}

\begin{proof}
We proceed by induction. The base case $t=1$ follows immediately from Equation~\eqref{eq:forward_step} with $\bar{\alpha}_1 = \alpha_1 = 1 - \beta_1$.

For the inductive step, assume the result holds for $t-1$. Then:
\begin{align*}
q(\mathbf{x}_t | \mathbf{x}_0) &= \int q(\mathbf{x}_t | \mathbf{x}_{t-1}) q(\mathbf{x}_{t-1} | \mathbf{x}_0) \, d\mathbf{x}_{t-1} \\
&= \int \N(\mathbf{x}_t; \sqrt{\alpha_t} \mathbf{x}_{t-1}, \beta_t \mathbf{I}) \N(\mathbf{x}_{t-1}; \sqrt{\bar{\alpha}_{t-1}} \mathbf{x}_0, (1 - \bar{\alpha}_{t-1}) \mathbf{I}) \, d\mathbf{x}_{t-1}.
\end{align*}

This is a convolution of two Gaussians. Using the identity for Gaussian convolution:
\[
\N(\mathbf{y}; A\mathbf{x}, \Sigma_1) * \N(\mathbf{x}; \boldsymbol{\mu}, \Sigma_2) = \N(\mathbf{y}; A\boldsymbol{\mu}, A\Sigma_2 A^\top + \Sigma_1),
\]

we obtain:
\begin{align*}
q(\mathbf{x}_t | \mathbf{x}_0) &= \N\left(\mathbf{x}_t; \sqrt{\alpha_t} \sqrt{\bar{\alpha}_{t-1}} \mathbf{x}_0, \alpha_t (1 - \bar{\alpha}_{t-1}) \mathbf{I} + \beta_t \mathbf{I}\right) \\
&= \N\left(\mathbf{x}_t; \sqrt{\alpha_t \bar{\alpha}_{t-1}} \mathbf{x}_0, (\alpha_t - \alpha_t \bar{\alpha}_{t-1} + \beta_t) \mathbf{I}\right) \\
&= \N\left(\mathbf{x}_t; \sqrt{\bar{\alpha}_t} \mathbf{x}_0, (1 - \bar{\alpha}_t) \mathbf{I}\right),
\end{align*}

where we used $\alpha_t \bar{\alpha}_{t-1} = \bar{\alpha}_t$ and $\alpha_t - \alpha_t \bar{\alpha}_{t-1} + \beta_t = \alpha_t(1 - \bar{\alpha}_{t-1}) + (1 - \alpha_t) = 1 - \alpha_t \bar{\alpha}_{t-1} = 1 - \bar{\alpha}_t$.
\end{proof}

\subsubsection{Direct Sampling via Reparameterization}

Equation~\eqref{eq:forward_closed} enables direct sampling of $\mathbf{x}_t$ given $\mathbf{x}_0$ without iterating through intermediate timesteps:

\begin{equation}
\mathbf{x}_t = \sqrt{\bar{\alpha}_t} \mathbf{x}_0 + \sqrt{1 - \bar{\alpha}_t} \boldsymbol{\epsilon}, \quad \boldsymbol{\epsilon} \sim \N(\mathbf{0}, \mathbf{I}).
\label{eq:xt_reparam}
\end{equation}

This formulation explicitly separates the signal contribution $\sqrt{\bar{\alpha}_t} \mathbf{x}_0$ from the noise contribution $\sqrt{1 - \bar{\alpha}_t} \boldsymbol{\epsilon}$.

\subsubsection{Variance Schedule Design}

The choice of variance schedule $\{\beta_t\}_{t=1}^T$ significantly impacts model performance. Common schedules include:

\begin{description}
\item[Linear Schedule] \citep{ho2020denoising}:
\begin{equation}
\beta_t = \beta_{\min} + \frac{t-1}{T-1}(\beta_{\max} - \beta_{\min}),
\end{equation}
typically with $\beta_{\min} = 10^{-4}$ and $\beta_{\max} = 0.02$.

\item[Cosine Schedule] \citep{nichol2021improved}:
\begin{equation}
\bar{\alpha}_t = \frac{f(t)}{f(0)}, \quad f(t) = \cos^2\left(\frac{t/T + s}{1 + s} \cdot \frac{\pi}{2}\right),
\end{equation}
where $s$ is a small offset (e.g., $s = 0.008$).

\item[Learned Schedule]: Treating $\{\beta_t\}$ or $\{\bar{\alpha}_t\}$ as learnable parameters optimized jointly with model weights.
\end{description}

\subsection{Reverse Diffusion Process}

The generative process involves reversing the forward diffusion. The true reverse conditional $q(\mathbf{x}_{t-1} | \mathbf{x}_t)$ is intractable as it requires integrating over all possible paths. However, for sufficiently small $\beta_t$, the reverse process is also Gaussian~\citep{feller1949theory, anderson1982reverse}.

We approximate the reverse process using a parameterized model $p_\theta$:

\begin{equation}
p_\theta(\mathbf{x}_{t-1} | \mathbf{x}_t) = \N(\mathbf{x}_{t-1}; \boldsymbol{\mu}_\theta(\mathbf{x}_t, t), \Sigma_\theta(\mathbf{x}_t, t)),
\label{eq:reverse_step}
\end{equation}

where $\boldsymbol{\mu}_\theta$ and $\Sigma_\theta$ are neural network outputs. For simplicity, the covariance is often fixed to $\Sigma_\theta(\mathbf{x}_t, t) = \sigma_t^2 \mathbf{I}$, where $\sigma_t$ is chosen according to:

\begin{equation}
\sigma_t^2 = \beta_t \quad \text{or} \quad \sigma_t^2 = \tilde{\beta}_t = \frac{1 - \bar{\alpha}_{t-1}}{1 - \bar{\alpha}_t} \beta_t.
\label{eq:variance_choices}
\end{equation}

\subsubsection{Posterior Distribution}

A crucial insight is that the true reverse conditional can be computed analytically when conditioning on $\mathbf{x}_0$:

\begin{lemma}[Reverse Posterior]
\begin{equation}
q(\mathbf{x}_{t-1} | \mathbf{x}_t, \mathbf{x}_0) = \N(\mathbf{x}_{t-1}; \tilde{\boldsymbol{\mu}}_t(\mathbf{x}_t, \mathbf{x}_0), \tilde{\beta}_t \mathbf{I}),
\label{eq:reverse_posterior}
\end{equation}
where
\begin{equation}
\tilde{\boldsymbol{\mu}}_t(\mathbf{x}_t, \mathbf{x}_0) = \frac{\sqrt{\bar{\alpha}_{t-1}} \beta_t}{1 - \bar{\alpha}_t} \mathbf{x}_0 + \frac{\sqrt{\alpha_t}(1 - \bar{\alpha}_{t-1})}{1 - \bar{\alpha}_t} \mathbf{x}_t.
\label{eq:posterior_mean}
\end{equation}
\end{lemma}

\begin{proof}
By Bayes' theorem:
\[
q(\mathbf{x}_{t-1} | \mathbf{x}_t, \mathbf{x}_0) = \frac{q(\mathbf{x}_t | \mathbf{x}_{t-1}, \mathbf{x}_0) q(\mathbf{x}_{t-1} | \mathbf{x}_0)}{q(\mathbf{x}_t | \mathbf{x}_0)}.
\]

By the Markov property, $q(\mathbf{x}_t | \mathbf{x}_{t-1}, \mathbf{x}_0) = q(\mathbf{x}_t | \mathbf{x}_{t-1})$. Substituting the Gaussian forms:

\begin{align*}
q(\mathbf{x}_{t-1} | \mathbf{x}_t, \mathbf{x}_0) &\propto \exp\left(-\frac{1}{2}\left[\frac{(\mathbf{x}_t - \sqrt{\alpha_t}\mathbf{x}_{t-1})^2}{\beta_t} + \frac{(\mathbf{x}_{t-1} - \sqrt{\bar{\alpha}_{t-1}}\mathbf{x}_0)^2}{1 - \bar{\alpha}_{t-1}} - \frac{(\mathbf{x}_t - \sqrt{\bar{\alpha}_t}\mathbf{x}_0)^2}{1 - \bar{\alpha}_t}\right]\right).
\end{align*}

Expanding and completing the square in $\mathbf{x}_{t-1}$:

\begin{align*}
&\frac{1}{\beta_t}(\mathbf{x}_t - \sqrt{\alpha_t}\mathbf{x}_{t-1})^2 + \frac{1}{1 - \bar{\alpha}_{t-1}}(\mathbf{x}_{t-1} - \sqrt{\bar{\alpha}_{t-1}}\mathbf{x}_0)^2 \\
&= \left(\frac{\alpha_t}{\beta_t} + \frac{1}{1 - \bar{\alpha}_{t-1}}\right)\mathbf{x}_{t-1}^2 - 2\left(\frac{\sqrt{\alpha_t}}{\beta_t}\mathbf{x}_t + \frac{\sqrt{\bar{\alpha}_{t-1}}}{1 - \bar{\alpha}_{t-1}}\mathbf{x}_0\right)\mathbf{x}_{t-1} + \text{const}.
\end{align*}

The coefficient of $\mathbf{x}_{t-1}^2$ gives the precision:
\[
\tilde{\beta}_t^{-1} = \frac{\alpha_t}{\beta_t} + \frac{1}{1 - \bar{\alpha}_{t-1}} = \frac{\alpha_t(1-\bar{\alpha}_{t-1}) + \beta_t}{\beta_t(1-\bar{\alpha}_{t-1})} = \frac{1 - \bar{\alpha}_t}{\beta_t(1-\bar{\alpha}_{t-1})},
\]

where we used $\alpha_t(1-\bar{\alpha}_{t-1}) + \beta_t = 1 - \bar{\alpha}_t$.

The mean follows from:
\[
\tilde{\boldsymbol{\mu}}_t = \tilde{\beta}_t\left(\frac{\sqrt{\alpha_t}}{\beta_t}\mathbf{x}_t + \frac{\sqrt{\bar{\alpha}_{t-1}}}{1 - \bar{\alpha}_{t-1}}\mathbf{x}_0\right).
\]

Substituting $\tilde{\beta}_t = \frac{(1-\bar{\alpha}_{t-1})\beta_t}{1-\bar{\alpha}_t}$ yields Equation~\eqref{eq:posterior_mean}.
\end{proof}

\subsection{Training Objective: Variational Lower Bound}

The training objective is derived from maximizing the log-likelihood $\log p_\theta(\mathbf{x}_0)$ of the data. We use the variational lower bound (VLB):

\begin{theorem}[Variational Lower Bound for Diffusion Models]
\begin{equation}
\log p_\theta(\mathbf{x}_0) \geq \E_{q(\mathbf{x}_{1:T}|\mathbf{x}_0)}\left[\log \frac{p_\theta(\mathbf{x}_{0:T})}{q(\mathbf{x}_{1:T}|\mathbf{x}_0)}\right] = -\loss_{\text{VLB}},
\label{eq:vlb}
\end{equation}
where
\begin{equation}
\loss_{\text{VLB}} = \loss_T + \sum_{t=2}^T \loss_{t-1} + \loss_0,
\label{eq:vlb_decomposed}
\end{equation}
with
\begin{align}
\loss_T &= \KL(q(\mathbf{x}_T|\mathbf{x}_0) \| p(\mathbf{x}_T)), \label{eq:loss_T} \\
\loss_{t-1} &= \KL(q(\mathbf{x}_{t-1}|\mathbf{x}_t, \mathbf{x}_0) \| p_\theta(\mathbf{x}_{t-1}|\mathbf{x}_t)), \label{eq:loss_t} \\
\loss_0 &= -\log p_\theta(\mathbf{x}_0|\mathbf{x}_1). \label{eq:loss_0}
\end{align}
\end{theorem}

\begin{proof}
Starting from the log-likelihood and introducing the forward process $q$:
\begin{align*}
\log p_\theta(\mathbf{x}_0) &= \log \int p_\theta(\mathbf{x}_{0:T}) \, d\mathbf{x}_{1:T} \\
&= \log \int \frac{p_\theta(\mathbf{x}_{0:T}) q(\mathbf{x}_{1:T}|\mathbf{x}_0)}{q(\mathbf{x}_{1:T}|\mathbf{x}_0)} \, d\mathbf{x}_{1:T} \\
&= \log \E_{q(\mathbf{x}_{1:T}|\mathbf{x}_0)}\left[\frac{p_\theta(\mathbf{x}_{0:T})}{q(\mathbf{x}_{1:T}|\mathbf{x}_0)}\right] \\
&\geq \E_{q(\mathbf{x}_{1:T}|\mathbf{x}_0)}\left[\log \frac{p_\theta(\mathbf{x}_{0:T})}{q(\mathbf{x}_{1:T}|\mathbf{x}_0)}\right] \quad \text{(Jensen's inequality)}.
\end{align*}

Expanding the joint distributions:
\begin{align*}
p_\theta(\mathbf{x}_{0:T}) &= p(\mathbf{x}_T) \prod_{t=1}^T p_\theta(\mathbf{x}_{t-1}|\mathbf{x}_t), \\
q(\mathbf{x}_{1:T}|\mathbf{x}_0) &= \prod_{t=1}^T q(\mathbf{x}_t|\mathbf{x}_{t-1}).
\end{align*}

Therefore:
\begin{align*}
\E_q\left[\log \frac{p_\theta(\mathbf{x}_{0:T})}{q(\mathbf{x}_{1:T}|\mathbf{x}_0)}\right] &= \E_q\left[\log p(\mathbf{x}_T) + \sum_{t=1}^T \log p_\theta(\mathbf{x}_{t-1}|\mathbf{x}_t) - \sum_{t=1}^T \log q(\mathbf{x}_t|\mathbf{x}_{t-1})\right].
\end{align*}

Using $q(\mathbf{x}_t|\mathbf{x}_{t-1}) = q(\mathbf{x}_t|\mathbf{x}_{t-1}, \mathbf{x}_0)$ and telescoping:
\begin{align*}
\sum_{t=1}^T \log q(\mathbf{x}_t|\mathbf{x}_{t-1}) &= \log q(\mathbf{x}_T|\mathbf{x}_0) + \sum_{t=1}^{T-1} \log \frac{q(\mathbf{x}_t|\mathbf{x}_{t-1}, \mathbf{x}_0) q(\mathbf{x}_t|\mathbf{x}_0)}{q(\mathbf{x}_{t+1}|\mathbf{x}_0)} \\
&= \log q(\mathbf{x}_T|\mathbf{x}_0) + \sum_{t=2}^T \log q(\mathbf{x}_{t-1}|\mathbf{x}_t, \mathbf{x}_0) - \log q(\mathbf{x}_1|\mathbf{x}_0).
\end{align*}

Collecting terms yields Equations~\eqref{eq:loss_T}--\eqref{eq:loss_0}.
\end{proof}

\subsubsection{Simplified Training Objective}

The term $\loss_T$ is constant (no trainable parameters), and $\loss_0$ can be modeled with a separate decoder. The key training signal comes from $\loss_{t-1}$. Since both $q$ and $p_\theta$ are Gaussians with fixed variance, the KL divergence simplifies to:

\begin{equation}
\loss_{t-1} = \E_{\mathbf{x}_0, \boldsymbol{\epsilon}}\left[\frac{1}{2\sigma_t^2}\norm{\tilde{\boldsymbol{\mu}}_t(\mathbf{x}_t, \mathbf{x}_0) - \boldsymbol{\mu}_\theta(\mathbf{x}_t, t)}^2\right].
\label{eq:loss_t_mean}
\end{equation}

\subsection{Noise Prediction Parameterization}

Rather than directly predicting $\boldsymbol{\mu}_\theta$, modern implementations predict the noise $\boldsymbol{\epsilon}$ added during the forward process~\citep{ho2020denoising}. Using Equation~\eqref{eq:xt_reparam}:

\begin{equation}
\mathbf{x}_0 = \frac{1}{\sqrt{\bar{\alpha}_t}}(\mathbf{x}_t - \sqrt{1 - \bar{\alpha}_t}\boldsymbol{\epsilon}).
\label{eq:x0_from_noise}
\end{equation}

Substituting into Equation~\eqref{eq:posterior_mean}:

\begin{equation}
\tilde{\boldsymbol{\mu}}_t(\mathbf{x}_t, \mathbf{x}_0) = \frac{1}{\sqrt{\alpha_t}}\left(\mathbf{x}_t - \frac{\beta_t}{\sqrt{1 - \bar{\alpha}_t}}\boldsymbol{\epsilon}\right).
\label{eq:mu_from_noise}
\end{equation}

Thus, we parameterize:

\begin{equation}
\boldsymbol{\mu}_\theta(\mathbf{x}_t, t) = \frac{1}{\sqrt{\alpha_t}}\left(\mathbf{x}_t - \frac{\beta_t}{\sqrt{1 - \bar{\alpha}_t}}\boldsymbol{\epsilon}_\theta(\mathbf{x}_t, t)\right),
\label{eq:mu_theta_noise}
\end{equation}

where $\boldsymbol{\epsilon}_\theta: \R^d \times \{1, \ldots, T\} \to \R^d$ is a neural network predicting the noise.

The simplified training objective becomes:

\begin{equation}
\loss_{\text{simple}} = \E_{t \sim \mathcal{U}(1,T), \mathbf{x}_0, \boldsymbol{\epsilon}}\left[\norm{\boldsymbol{\epsilon} - \boldsymbol{\epsilon}_\theta(\mathbf{x}_t, t)}^2\right],
\label{eq:loss_simple}
\end{equation}

where $\mathbf{x}_t = \sqrt{\bar{\alpha}_t}\mathbf{x}_0 + \sqrt{1-\bar{\alpha}_t}\boldsymbol{\epsilon}$ and $\boldsymbol{\epsilon} \sim \N(\mathbf{0}, \mathbf{I})$.

\begin{remark}
Ho et al.~\citep{ho2020denoising} found empirically that dropping the weighting factor $\frac{1}{2\sigma_t^2}$ improves sample quality, though this deviates from the strict VLB objective.
\end{remark}

\section{Score-Based Perspective}

An alternative but equivalent formulation views diffusion models as learning the score function $\nabla_{\mathbf{x}} \log q(\mathbf{x})$~\citep{song2019generative, song2020score}.

\subsection{Score Function and Langevin Dynamics}

The score function measures the gradient of the log-density. Given a smooth density $p(\mathbf{x})$, samples can be generated via Langevin dynamics:

\begin{equation}
\mathbf{x}_{k+1} = \mathbf{x}_k + \frac{\eta}{2}\nabla_{\mathbf{x}} \log p(\mathbf{x}_k) + \sqrt{\eta}\mathbf{z}_k, \quad \mathbf{z}_k \sim \N(\mathbf{0}, \mathbf{I}),
\label{eq:langevin}
\end{equation}

which converges to samples from $p(\mathbf{x})$ as $\eta \to 0$ and $k \to \infty$.

\subsection{Denoising Score Matching}

Direct score matching is challenging due to the need to compute $\nabla_{\mathbf{x}} \log q(\mathbf{x})$. Vincent~\citep{vincent2011connection} showed that the score can be learned by training a model to denoise corrupted data:

\begin{theorem}[Denoising Score Matching]
Let $q_\sigma(\tilde{\mathbf{x}}|\mathbf{x}) = \N(\tilde{\mathbf{x}}; \mathbf{x}, \sigma^2\mathbf{I})$ be a noise distribution. Then:
\begin{equation}
\nabla_{\mathbf{x}} \log q_\sigma(\mathbf{x}) = \E_{q_\sigma(\tilde{\mathbf{x}}|\mathbf{x})}\left[\nabla_{\mathbf{x}} \log q_\sigma(\tilde{\mathbf{x}}|\mathbf{x})\right] = \E_{\boldsymbol{\epsilon}}\left[-\frac{\boldsymbol{\epsilon}}{\sigma}\right] = -\frac{1}{\sigma^2}\E_{q_\sigma(\tilde{\mathbf{x}}|\mathbf{x})}[\tilde{\mathbf{x}} - \mathbf{x}].
\label{eq:score_denoising}
\end{equation}
\end{theorem}

Thus, training a model $\mathbf{s}_\theta(\mathbf{x}, \sigma)$ to predict $\frac{\mathbf{x} - \tilde{\mathbf{x}}}{\sigma^2}$ is equivalent to estimating the score.

\subsection{Connection to Diffusion Models}

For diffusion models, the score of the forward marginal is:

\begin{equation}
\nabla_{\mathbf{x}_t} \log q(\mathbf{x}_t|\mathbf{x}_0) = -\frac{1}{1 - \bar{\alpha}_t}(\mathbf{x}_t - \sqrt{\bar{\alpha}_t}\mathbf{x}_0) = -\frac{1}{\sqrt{1 - \bar{\alpha}_t}}\boldsymbol{\epsilon}.
\label{eq:score_diffusion}
\end{equation}

Therefore, predicting the noise $\boldsymbol{\epsilon}_\theta(\mathbf{x}_t, t)$ is equivalent to predicting the score (up to scaling):

\begin{equation}
\mathbf{s}_\theta(\mathbf{x}_t, t) = -\frac{1}{\sqrt{1 - \bar{\alpha}_t}}\boldsymbol{\epsilon}_\theta(\mathbf{x}_t, t) \approx \nabla_{\mathbf{x}_t} \log q(\mathbf{x}_t).
\label{eq:score_noise_relation}
\end{equation}

\section{Stochastic Differential Equation Formulation}

Song et al.~\citep{song2020score} unified diffusion models and score-based models through the lens of stochastic differential equations.

\subsection{Forward SDE}

The continuous-time limit of the forward diffusion process is described by the SDE:

\begin{equation}
d\mathbf{x} = \mathbf{f}(\mathbf{x}, t) \, dt + g(t) \, d\mathbf{w},
\label{eq:forward_sde}
\end{equation}

where $\mathbf{w}$ is a standard Wiener process, $\mathbf{f}(\mathbf{x}, t)$ is the drift coefficient, and $g(t)$ is the diffusion coefficient.

For the variance-preserving (VP) formulation:

\begin{equation}
d\mathbf{x} = -\frac{1}{2}\beta(t)\mathbf{x} \, dt + \sqrt{\beta(t)} \, d\mathbf{w}.
\label{eq:vp_sde}
\end{equation}

For the variance-exploding (VE) formulation:

\begin{equation}
d\mathbf{x} = \sqrt{\frac{d[\sigma^2(t)]}{dt}} \, d\mathbf{w}.
\label{eq:ve_sde}
\end{equation}

\subsection{Reverse SDE}

Anderson~\citep{anderson1982reverse} showed that the reverse-time SDE is:

\begin{equation}
d\mathbf{x} = \left[\mathbf{f}(\mathbf{x}, t) - g^2(t) \nabla_{\mathbf{x}} \log p_t(\mathbf{x})\right] dt + g(t) \, d\bar{\mathbf{w}},
\label{eq:reverse_sde}
\end{equation}

where $\bar{\mathbf{w}}$ is a reverse-time Wiener process and $p_t$ is the marginal density at time $t$.

Substituting the score approximation $\nabla_{\mathbf{x}} \log p_t(\mathbf{x}) \approx \mathbf{s}_\theta(\mathbf{x}, t)$:

\begin{equation}
d\mathbf{x} = \left[\mathbf{f}(\mathbf{x}, t) - g^2(t) \mathbf{s}_\theta(\mathbf{x}, t)\right] dt + g(t) \, d\bar{\mathbf{w}}.
\label{eq:reverse_sde_learned}
\end{equation}

This SDE can be integrated numerically using methods like Euler-Maruyama or higher-order solvers~\citep{jolicoeur2021gotta}.

\subsection{Probability Flow ODE}

Song et al.~\citep{song2020score} derived a deterministic ODE that generates the same marginal densities:

\begin{equation}
\frac{d\mathbf{x}}{dt} = \mathbf{f}(\mathbf{x}, t) - \frac{1}{2}g^2(t) \nabla_{\mathbf{x}} \log p_t(\mathbf{x}).
\label{eq:prob_flow_ode}
\end{equation}

This formulation enables:
\begin{itemize}
\item Exact likelihood computation via change of variables.
\item Latent space manipulation and interpolation.
\item Connection to neural ODEs and normalizing flows.
\end{itemize}

\section{Classifier-Free Guidance}

For conditional generation $p(\mathbf{x}|y)$ where $y$ is a class label or text prompt, classifier-free guidance~\citep{ho2022classifier} modifies the score:

\begin{equation}
\tilde{\boldsymbol{\epsilon}}_\theta(\mathbf{x}_t, t, y) = \boldsymbol{\epsilon}_\theta(\mathbf{x}_t, t, \emptyset) + s \cdot (\boldsymbol{\epsilon}_\theta(\mathbf{x}_t, t, y) - \boldsymbol{\epsilon}_\theta(\mathbf{x}_t, t, \emptyset)),
\label{eq:cfg}
\end{equation}

where $s \geq 1$ is the guidance scale and $\emptyset$ denotes the unconditional (null) class. This corresponds to sampling from:

\begin{equation}
\tilde{p}(\mathbf{x}_t) \propto p(\mathbf{x}_t)^{1-s} p(\mathbf{x}_t|y)^s.
\label{eq:cfg_distribution}
\end{equation}

Higher guidance scales increase sample fidelity to the condition but may reduce diversity.

\chapter{Model Compression Techniques: Theory and Practice}

This chapter presents comprehensive theoretical foundations for three primary model compression techniques applicable to diffusion models: knowledge distillation, structured pruning, and quantization. For each method, we provide complete mathematical derivations, convergence analyses, and practical considerations.

\section{Knowledge Distillation}

Knowledge distillation~\citep{hinton2015distilling} transfers knowledge from a large teacher model to a compact student model by training the student to mimic the teacher's behavior.

\subsection{Classical Distillation Framework}

\subsubsection{Temperature-Scaled Softmax}

For classification tasks, let $\mathbf{z}_T$ and $\mathbf{z}_S$ denote teacher and student logits. The temperature-scaled softmax is:

\begin{equation}
p_i^{(T)}(\mathbf{z}_T, \tau) = \frac{\exp(z_{T,i}/\tau)}{\sum_j \exp(z_{T,j}/\tau)},
\label{eq:temp_softmax}
\end{equation}

where $\tau > 0$ is the temperature. As $\tau \to 0$, the distribution becomes peaked (one-hot), while $\tau \to \infty$ yields a uniform distribution.

\subsubsection{Distillation Loss}

The total loss combines task loss and distillation loss:

\begin{equation}
\loss_{\text{distill}} = \alpha \loss_{\text{task}}(\mathbf{z}_S, y) + \beta \tau^2 \KL(p^{(T)}(\mathbf{z}_T, \tau) \| p^{(S)}(\mathbf{z}_S, \tau)),
\label{eq:distill_loss_classification}
\end{equation}

where $\alpha, \beta$ are weighting coefficients and the $\tau^2$ factor accounts for gradient scaling.

\subsection{Distillation for Diffusion Models}

For diffusion models, distillation operates on the noise prediction function $\boldsymbol{\epsilon}_\theta(\mathbf{x}_t, t)$.

\subsubsection{Direct Output Matching}

The simplest approach matches student predictions to teacher predictions:

\begin{equation}
\loss_{\text{distill}}^{\text{output}} = \E_{t, \mathbf{x}_0, \boldsymbol{\epsilon}}\left[\norm{\boldsymbol{\epsilon}_T(\mathbf{x}_t, t) - \boldsymbol{\epsilon}_S(\mathbf{x}_t, t)}^2\right],
\label{eq:output_distill}
\end{equation}

where subscripts $T$ and $S$ denote teacher and student.

\subsubsection{Combined Objective}

The student is trained with a weighted combination:

\begin{equation}
\loss_{\text{total}} = \alpha \loss_{\text{diffusion}} + \beta \loss_{\text{distill}}^{\text{output}},
\label{eq:combined_distill_loss}
\end{equation}

where $\loss_{\text{diffusion}}$ is the standard diffusion objective (Equation~\eqref{eq:loss_simple}).

\subsubsection{Intermediate Feature Matching}

Beyond final outputs, matching intermediate representations can improve distillation~\citep{romero2014fitnets}. Let $\mathbf{F}_T^{(\ell)}$ and $\mathbf{F}_S^{(\ell)}$ denote activations from layer $\ell$ of teacher and student. The feature matching loss is:

\begin{equation}
\loss_{\text{feature}}^{(\ell)} = \E\left[\norm{\mathbf{F}_T^{(\ell)}(\mathbf{x}_t, t) - \Phi^{(\ell)}(\mathbf{F}_S^{(\ell)}(\mathbf{x}_t, t))}^2\right],
\label{eq:feature_matching}
\end{equation}

where $\Phi^{(\ell)}$ is an optional projection layer to align dimensions.

The total distillation loss becomes:

\begin{equation}
\loss_{\text{distill}} = \loss_{\text{distill}}^{\text{output}} + \sum_{\ell \in \mathcal{L}} \gamma_\ell \loss_{\text{feature}}^{(\ell)},
\label{eq:total_distill_feature}
\end{equation}

where $\mathcal{L}$ is the set of layers for feature matching and $\gamma_\ell$ are layer-specific weights.

\subsection{Progressive Distillation for Diffusion}

Salimans and Ho~\citep{salimans2022progressive} proposed progressive distillation specifically for accelerating diffusion model sampling.

\subsubsection{Two-Step Distillation}

The key idea is to train the student to predict in one step what the teacher predicts in two steps. Let $\mathbf{x}_t^{(0)} = \mathbf{x}_t$ and define:

\begin{equation}
\mathbf{x}_t^{(1)} = \sqrt{\bar{\alpha}_{t-1}}\hat{\mathbf{x}}_0^{(0)} + \sqrt{1 - \bar{\alpha}_{t-1}}\boldsymbol{\epsilon}_T(\mathbf{x}_t^{(0)}, t),
\label{eq:x_t_1}
\end{equation}

where $\hat{\mathbf{x}}_0^{(0)} = \frac{1}{\sqrt{\bar{\alpha}_t}}(\mathbf{x}_t^{(0)} - \sqrt{1-\bar{\alpha}_t}\boldsymbol{\epsilon}_T(\mathbf{x}_t^{(0)}, t))$.

Then compute:

\begin{equation}
\mathbf{x}_t^{(2)} = \sqrt{\bar{\alpha}_{t-2}}\hat{\mathbf{x}}_0^{(1)} + \sqrt{1 - \bar{\alpha}_{t-2}}\boldsymbol{\epsilon}_T(\mathbf{x}_t^{(1)}, t-1).
\label{eq:x_t_2}
\end{equation}

The student is trained to predict directly from $\mathbf{x}_t^{(0)}$ to $\mathbf{x}_t^{(2)}$:

\begin{equation}
\loss_{\text{progressive}} = \E\left[\norm{\mathbf{x}_t^{(2)} - \left(\sqrt{\bar{\alpha}_{t-2}}\hat{\mathbf{x}}_0^{(S)} + \sqrt{1 - \bar{\alpha}_{t-2}}\boldsymbol{\epsilon}_S(\mathbf{x}_t, t)\right)}^2\right].
\label{eq:progressive_loss}
\end{equation}

\subsubsection{Iterative Halving}

By repeatedly applying this procedure, the number of sampling steps can be reduced exponentially: $T \to T/2 \to T/4 \to \cdots \to 1$.

\subsection{Consistency Distillation}

Song et al.~\citep{song2023consistency} introduced consistency models that map any point on a diffusion trajectory to the origin:

\begin{equation}
\mathbf{f}_\theta(\mathbf{x}_t, t) \approx \mathbf{x}_0.
\label{eq:consistency_map}
\end{equation}

The consistency loss enforces:

\begin{equation}
\loss_{\text{consistency}} = \E\left[\norm{\mathbf{f}_\theta(\mathbf{x}_{t+\Delta t}, t+\Delta t) - \mathbf{f}_\theta(\mathbf{x}_t, t)}^2\right],
\label{eq:consistency_loss}
\end{equation}

where $\mathbf{x}_{t+\Delta t}$ is obtained by one step of the diffusion ODE starting from $\mathbf{x}_t$.

This enables single-step generation: $\mathbf{x}_0 \approx \mathbf{f}_\theta(\mathbf{x}_T, T)$ where $\mathbf{x}_T \sim \N(\mathbf{0}, \mathbf{I})$.

\section{Structured Pruning}

Pruning reduces model complexity by removing parameters based on importance criteria~\citep{lecun1990optimal, han2015learning}.

\subsection{Pruning Formulations}

\subsubsection{Unstructured Pruning}

Unstructured pruning removes individual weights:

\begin{equation}
w_{ij} \leftarrow w_{ij} \cdot m_{ij}, \quad m_{ij} \in \{0, 1\},
\label{eq:unstructured_prune}
\end{equation}

where $m_{ij}$ is a binary mask. The mask is typically determined by magnitude thresholding:

\begin{equation}
m_{ij} = \mathbbm{1}_{|w_{ij}| \geq \delta_p},
\label{eq:magnitude_mask}
\end{equation}

where $\delta_p$ is chosen to achieve a target sparsity $s$ (fraction of pruned weights).

\subsubsection{Structured Pruning}

Structured pruning removes entire channels, neurons, or blocks:

\begin{equation}
\mathbf{W}_{:,j} \leftarrow \mathbf{W}_{:,j} \cdot m_j, \quad m_j \in \{0, 1\},
\label{eq:structured_prune}
\end{equation}

where $\mathbf{W}_{:,j}$ denotes the $j$-th output channel (column).

\subsection{Importance Criteria}

\subsubsection{Magnitude-Based Pruning}

The simplest criterion is weight magnitude:

\begin{equation}
\text{importance}(w_{ij}) = |w_{ij}|.
\label{eq:importance_magnitude}
\end{equation}

For structured pruning of channels:

\begin{equation}
\text{importance}(\mathbf{W}_{:,j}) = \norm{\mathbf{W}_{:,j}}_1 \quad \text{or} \quad \norm{\mathbf{W}_{:,j}}_2.
\label{eq:importance_channel}
\end{equation}

\subsubsection{Gradient-Based Pruning}

Gradient magnitude reflects sensitivity:

\begin{equation}
\text{importance}(w_{ij}) = \left|w_{ij} \cdot \frac{\partial \loss}{\partial w_{ij}}\right|.
\label{eq:importance_gradient}
\end{equation}

\subsubsection{Second-Order Methods}

The Optimal Brain Damage (OBD) criterion~\citep{lecun1990optimal} uses second-order information:

\begin{equation}
\text{importance}(w_{ij}) = \frac{1}{2} h_{ij} w_{ij}^2,
\label{eq:importance_obd}
\end{equation}

where $h_{ij} = \frac{\partial^2 \loss}{\partial w_{ij}^2}$ is the Hessian diagonal.

\subsection{Iterative Pruning with Fine-Tuning}

Aggressive one-shot pruning often degrades performance significantly. Iterative pruning alternates between pruning and fine-tuning:

\begin{algorithm}[H]
\caption{Iterative Magnitude Pruning}
\label{alg:iterative_pruning}
\begin{algorithmic}[1]
\STATE \textbf{Input:} Trained model $\theta_0$, target sparsity $s_{\text{target}}$, iterations $N$
\STATE \textbf{Output:} Pruned model $\theta_N$
\FOR{$i = 1$ to $N$}
    \STATE $s_i \leftarrow s_{\text{target}} \cdot \left(1 - (1 - i/N)^3\right)$ \COMMENT{Cubic schedule}
    \STATE Compute importance scores for all weights
    \STATE Set bottom $s_i$ fraction of weights to zero
    \STATE Fine-tune for $K$ steps to recover performance
\ENDFOR
\RETURN $\theta_N$
\end{algorithmic}
\end{algorithm}

\subsection{Dependency Graphs for Structured Pruning}

When pruning entire channels, dependencies must be tracked:

\begin{itemize}
\item \textbf{Conv Layer}: Pruning output channel $j$ removes $\mathbf{W}_{:,j,:,:}$ and requires removing corresponding input channel $j$ in the next layer.
\item \textbf{Skip Connections}: Residual connections constrain pruning—both branches must maintain compatible dimensions.
\item \textbf{Batch Normalization}: Pruning a channel requires removing corresponding BN parameters.
\end{itemize}

Formally, construct a dependency graph $\mathcal{G} = (\mathcal{V}, \mathcal{E})$ where:
\begin{itemize}
\item Vertices $\mathcal{V}$ represent channels/neurons.
\item Edges $(\mathbf{v}_i, \mathbf{v}_j) \in \mathcal{E}$ indicate that pruning $\mathbf{v}_i$ requires pruning $\mathbf{v}_j$.
\end{itemize}

Pruning must respect the dependency graph to maintain architectural integrity.

\subsection{Layer-Wise Sensitivity Analysis}

Not all layers contribute equally to performance. We define layer sensitivity as:

\begin{definition}[Layer Sensitivity]
The sensitivity of layer $\ell$ to pruning is:
\begin{equation}
S^{(\ell)}(p) = \loss_{\text{pruned}}^{(\ell)}(p) - \loss_{\text{original}},
\label{eq:layer_sensitivity}
\end{equation}
where $\loss_{\text{pruned}}^{(\ell)}(p)$ is the loss after pruning fraction $p$ of layer $\ell$'s parameters.
\end{definition}

Layers with lower sensitivity can tolerate higher pruning rates. This leads to non-uniform pruning strategies:

\begin{equation}
p^{(\ell)} = p_{\text{base}} \cdot \exp(-\lambda S^{(\ell)}(p_{\text{probe}})),
\label{eq:adaptive_pruning}
\end{equation}

where $p_{\text{base}}$ is a baseline rate, $\lambda$ controls sensitivity scaling, and $p_{\text{probe}}$ is a small probe pruning fraction for estimating $S^{(\ell)}$.

\subsection{Pruning for Diffusion Models}

Diffusion models present unique challenges:

\begin{enumerate}
\item \textbf{Temporal Consistency}: Pruning must maintain consistency across timesteps.
\item \textbf{Multi-Resolution Features}: U-Net architectures have skip connections across resolutions.
\item \textbf{Attention Layers}: Self-attention is parameter-efficient but computationally expensive—pruning attention heads requires careful analysis.
\end{enumerate}

\subsubsection{Timestep-Aware Pruning}

Define per-timestep importance:

\begin{equation}
\text{importance}^{(t)}(\mathbf{W}_{:,j}) = \E_{\mathbf{x}_0, \boldsymbol{\epsilon}}\left[\left\|\frac{\partial \loss(\mathbf{x}_t, t)}{\partial \mathbf{W}_{:,j}}\right\|_2\right],
\label{eq:timestep_importance}
\end{equation}

and aggregate across critical timesteps:

\begin{equation}
\text{importance}(\mathbf{W}_{:,j}) = \sum_{t \in \mathcal{T}_{\text{critical}}} w_t \cdot \text{importance}^{(t)}(\mathbf{W}_{:,j}),
\label{eq:importance_aggregated}
\end{equation}

where $\mathcal{T}_{\text{critical}}$ contains timesteps contributing most to generation quality.

\section{Quantization}

Quantization reduces numerical precision of weights and activations~\citep{jacob2018quantization}.

\subsection{Uniform Quantization}

\subsubsection{Quantization Mapping}

Map continuous values to discrete levels:

\begin{equation}
w_q = \Delta \cdot \left\lfloor \frac{w - z}{\Delta} \right\rceil + z,
\label{eq:uniform_quant}
\end{equation}

where:
\begin{itemize}
\item $\Delta = \frac{w_{\max} - w_{\min}}{2^b - 1}$ is the quantization step (scale),
\item $b$ is the bit-width,
\item $z$ is the zero-point: $z = -\left\lfloor \frac{w_{\min}}{\Delta} \right\rceil \Delta$.
\end{itemize}

\subsubsection{Integer Representation}

For efficient computation, represent as integers:

\begin{equation}
w_q = \Delta(q_{\text{int}} - z_{\text{int}}),
\label{eq:int_representation}
\end{equation}

where $q_{\text{int}} \in \{0, 1, \ldots, 2^b-1\}$ and $z_{\text{int}} = \left\lfloor \frac{w_{\min}}{\Delta} \right\rceil$.

\subsection{Quantization-Aware Training (QAT)}

Post-training quantization can cause significant accuracy drop. QAT simulates quantization during training.

\subsubsection{Fake Quantization}

Insert quantization operations in forward pass:

\begin{equation}
\tilde{w} = Q(w) = \Delta \cdot \text{clip}\left(\left\lfloor \frac{w}{\Delta} \right\rceil, 0, 2^b - 1\right),
\label{eq:fake_quant}
\end{equation}

where $\text{clip}$ ensures values stay within representable range.

\subsubsection{Straight-Through Estimator (STE)}

The rounding operation is non-differentiable. The straight-through estimator~\citep{bengio2013estimating} approximates gradients:

\begin{equation}
\frac{\partial Q(w)}{\partial w} \approx \begin{cases}
1, & \text{if } w_{\min} \leq w \leq w_{\max}, \\
0, & \text{otherwise}.
\end{cases}
\label{eq:ste}
\end{equation}

This allows gradient backpropagation through quantization layers.

\subsection{Per-Channel vs Per-Tensor Quantization}

\subsubsection{Per-Tensor Quantization}

A single scale $\Delta$ for the entire tensor:

\begin{equation}
\Delta_{\text{tensor}} = \frac{\max(\mathbf{W}) - \min(\mathbf{W})}{2^b - 1}.
\label{eq:per_tensor_scale}
\end{equation}

Simple but suboptimal if different channels have different dynamic ranges.

\subsubsection{Per-Channel Quantization}

Separate scale for each output channel:

\begin{equation}
\Delta_j = \frac{\max(\mathbf{W}_{:,j}) - \min(\mathbf{W}_{:,j})}{2^b - 1}.
\label{eq:per_channel_scale}
\end{equation}

More accurate representation at the cost of storing multiple scales.

\subsection{Mixed-Precision Quantization}

Different layers may tolerate different bit-widths. Define bit-width allocation $\mathbf{b} = (b_1, \ldots, b_L)$ for $L$ layers. The optimization problem is:

\begin{equation}
\begin{aligned}
\minimize_{\mathbf{b}} \quad & \sum_{\ell=1}^L \text{size}^{(\ell)}(b_\ell) \\
\text{subject to} \quad & \loss_{\text{quant}}(\mathbf{b}) \leq \loss_{\text{target}}, \\
& b_\ell \in \{2, 4, 8, 16, 32\}, \quad \forall \ell.
\end{aligned}
\label{eq:mixed_precision_opt}
\end{equation}

This is typically solved via reinforcement learning or evolutionary algorithms~\citep{wang2019haq}.

\subsection{Quantization for Diffusion Models}

\subsubsection{Activation Quantization Challenges}

Activation distributions vary across timesteps in diffusion models. Using fixed quantization ranges can clip important values at certain timesteps.

\textbf{Solution}: Dynamic quantization with timestep-dependent ranges:

\begin{equation}
\Delta^{(t)} = f_\theta(t), \quad z^{(t)} = g_\theta(t),
\label{eq:dynamic_quant}
\end{equation}

where $f_\theta$ and $g_\theta$ are small learned networks predicting scale and zero-point.

\subsubsection{Quantization-Aware Distillation}

Combine quantization and distillation:

\begin{equation}
\loss = \alpha \loss_{\text{diffusion}} + \beta \loss_{\text{distill}} + \gamma \loss_{\text{quant}},
\label{eq:quant_distill_loss}
\end{equation}

where $\loss_{\text{quant}}$ encourages quantization-friendly weights:

\begin{equation}
\loss_{\text{quant}} = \E\left[\norm{w - Q(w)}^2\right].
\label{eq:quant_regularizer}
\end{equation}

\chapter{Efficient Training Techniques}

Training large diffusion models requires careful optimization of computational resources. This chapter covers mixed precision training, exponential moving averages, and gradient scaling strategies.

\section{Mixed Precision Training}

Mixed precision training~\citep{micikevicius2017mixed} combines FP16 and FP32 arithmetic to accelerate training while maintaining numerical stability.

\subsection{Floating-Point Representations}

\subsubsection{IEEE 754 Formats}

\begin{itemize}
\item \textbf{FP32 (Single Precision)}: 1 sign bit, 8 exponent bits, 23 mantissa bits.
  \begin{itemize}
  \item Range: $\approx 10^{-38}$ to $10^{38}$
  \item Precision: $\approx 7$ decimal digits
  \end{itemize}

\item \textbf{FP16 (Half Precision)}: 1 sign bit, 5 exponent bits, 10 mantissa bits.
  \begin{itemize}
  \item Range: $\approx 10^{-8}$ to $65504$
  \item Precision: $\approx 3$ decimal digits
  \end{itemize}
\end{itemize}

\subsection{Mixed Precision Training Algorithm}

\begin{algorithm}[H]
\caption{Mixed Precision Training with Loss Scaling}
\label{alg:mixed_precision}
\begin{algorithmic}[1]
\STATE \textbf{Input:} Model $\theta$ (FP32), loss scale $s$, data $\mathcal{D}$
\STATE \textbf{Output:} Trained model $\theta$
\WHILE{not converged}
    \STATE Sample batch $(\mathbf{x}, \mathbf{y}) \sim \mathcal{D}$
    \STATE Cast model to FP16: $\theta_{\text{FP16}} \leftarrow \text{cast}(\theta)$
    \STATE Forward pass in FP16: $\hat{\mathbf{y}} = f_{\theta_{\text{FP16}}}(\mathbf{x})$
    \STATE Compute loss in FP32: $\loss = \text{loss}(\hat{\mathbf{y}}, \mathbf{y})$
    \STATE Scale loss: $\loss_{\text{scaled}} = s \cdot \loss$
    \STATE Backward pass: compute gradients $\mathbf{g}_{\text{FP16}} = \nabla_\theta \loss_{\text{scaled}}$
    \STATE Unscale gradients: $\mathbf{g} = \mathbf{g}_{\text{FP16}} / s$
    \STATE Check for overflow/underflow; adjust $s$ if necessary
    \STATE Update FP32 master weights: $\theta \leftarrow \theta - \eta \mathbf{g}$
\ENDWHILE
\end{algorithmic}
\end{algorithm}

\subsection{Loss Scaling}

\subsubsection{Static Loss Scaling}

Use a fixed scale factor $s$ (e.g., $s = 2^{10} = 1024$):

\begin{equation}
\loss_{\text{scaled}} = s \cdot \loss.
\label{eq:static_loss_scale}
\end{equation}

Gradients are computed with respect to $\loss_{\text{scaled}}$ then divided by $s$:

\begin{equation}
\mathbf{g} = \frac{1}{s} \nabla_\theta \loss_{\text{scaled}} = \nabla_\theta \loss.
\label{eq:unscale_gradient}
\end{equation}

\subsubsection{Dynamic Loss Scaling}

Automatically adjust $s$ based on gradient statistics:

\begin{algorithm}[H]
\caption{Dynamic Loss Scaling}
\label{alg:dynamic_loss_scale}
\begin{algorithmic}[1]
\STATE Initialize $s = 2^{15}$, $s_{\min} = 1$, $s_{\max} = 2^{20}$
\STATE $n_{\text{stable}} = 0$, $n_{\text{threshold}} = 2000$
\WHILE{training}
    \IF{gradients contain NaN or Inf}
        \STATE $s \leftarrow \max(s / 2, s_{\min})$
        \STATE $n_{\text{stable}} \leftarrow 0$
        \STATE Skip weight update
    \ELSE
        \STATE $n_{\text{stable}} \leftarrow n_{\text{stable}} + 1$
        \IF{$n_{\text{stable}} \geq n_{\text{threshold}}$}
            \STATE $s \leftarrow \min(2s, s_{\max})$
            \STATE $n_{\text{stable}} \leftarrow 0$
        \ENDIF
    \ENDIF
\ENDWHILE
\end{algorithmic}
\end{algorithm}

\subsection{Selective Precision}

Certain operations must remain in FP32 for stability:

\begin{itemize}
\item \textbf{Batch Normalization}: Running mean/variance updated in FP32.
\item \textbf{Softmax}: Computed in FP32 to avoid overflow.
\item \textbf{Loss Computation}: Always FP32.
\item \textbf{Master Weights}: Stored in FP32; cast to FP16 for forward/backward.
\end{itemize}

\section{Exponential Moving Average (EMA)}

EMA smooths parameter updates, leading to better generalization~\citep{polyak1992acceleration}.

\subsection{EMA Update Rule}

Maintain shadow parameters $\theta_{\text{EMA}}$:

\begin{equation}
\theta_{\text{EMA}}^{(t+1)} = \beta \theta_{\text{EMA}}^{(t)} + (1 - \beta) \theta^{(t+1)},
\label{eq:ema_update}
\end{equation}

where $\beta \in (0, 1)$ is the decay rate (typically $\beta = 0.995$ or $\beta = 0.9999$).

\subsection{Bias Correction}

During early training, EMA is biased toward initialization. Apply bias correction:

\begin{equation}
\hat{\theta}_{\text{EMA}}^{(t)} = \frac{\theta_{\text{EMA}}^{(t)}}{1 - \beta^t}.
\label{eq:ema_bias_correction}
\end{equation}

\subsection{EMA for Batch Normalization}

Batch normalization layers maintain running statistics. Apply EMA separately:

\begin{align}
\mu_{\text{running}}^{(t+1)} &= \beta_{\text{BN}} \mu_{\text{running}}^{(t)} + (1 - \beta_{\text{BN}}) \mu_{\text{batch}}^{(t)}, \label{eq:bn_mean_ema} \\
\sigma_{\text{running}}^{2(t+1)} &= \beta_{\text{BN}} \sigma_{\text{running}}^{2(t)} + (1 - \beta_{\text{BN}}) \sigma_{\text{batch}}^{2(t)}, \label{eq:bn_var_ema}
\end{align}

where $\beta_{\text{BN}} \approx 0.9$ (faster decay than model weights).

\subsection{EMA Benefits for Diffusion Models}

\begin{enumerate}
\item \textbf{Improved Sample Quality}: EMA parameters often produce higher-quality samples than raw parameters.
\item \textbf{Training Stability}: Smoothing reduces variance from stochastic gradient noise.
\item \textbf{Reduced Overfitting}: EMA acts as implicit regularization.
\end{enumerate}

\section{Gradient Accumulation}

When memory constraints limit batch size, gradient accumulation simulates larger batches:

\begin{algorithm}[H]
\caption{Gradient Accumulation}
\label{alg:grad_accum}
\begin{algorithmic}[1]
\STATE \textbf{Input:} Model $\theta$, micro-batch size $B_{\text{micro}}$, accumulation steps $K$
\STATE \textbf{Output:} Updated model $\theta$
\STATE Initialize $\mathbf{g}_{\text{accum}} = \mathbf{0}$
\FOR{$k = 1$ to $K$}
    \STATE Sample micro-batch $\mathcal{B}_k$ of size $B_{\text{micro}}$
    \STATE Compute loss $\loss_k$ on $\mathcal{B}_k$
    \STATE Compute gradients $\mathbf{g}_k = \nabla_\theta \loss_k$
    \STATE $\mathbf{g}_{\text{accum}} \leftarrow \mathbf{g}_{\text{accum}} + \mathbf{g}_k$
\ENDFOR
\STATE $\mathbf{g}_{\text{accum}} \leftarrow \mathbf{g}_{\text{accum}} / K$
\STATE Update $\theta \leftarrow \theta - \eta \mathbf{g}_{\text{accum}}$
\end{algorithmic}
\end{algorithm}

Effective batch size: $B_{\text{eff}} = K \cdot B_{\text{micro}}$.

\section{Optimizer Selection}

\subsection{AdamW}

AdamW~\citep{loshchilov2017decoupled} decouples weight decay from gradient-based updates:

\begin{align}
\mathbf{m}_t &= \beta_1 \mathbf{m}_{t-1} + (1 - \beta_1) \mathbf{g}_t, \label{eq:adamw_m} \\
\mathbf{v}_t &= \beta_2 \mathbf{v}_{t-1} + (1 - \beta_2) \mathbf{g}_t^2, \label{eq:adamw_v} \\
\hat{\mathbf{m}}_t &= \mathbf{m}_t / (1 - \beta_1^t), \quad \hat{\mathbf{v}}_t = \mathbf{v}_t / (1 - \beta_2^t), \label{eq:adamw_bias_correct} \\
\theta_t &= \theta_{t-1} - \eta \left(\frac{\hat{\mathbf{m}}_t}{\sqrt{\hat{\mathbf{v}}_t} + \epsilon} + \lambda \theta_{t-1}\right), \label{eq:adamw_update}
\end{align}

where $\lambda$ is the weight decay coefficient.

\subsection{Learning Rate Schedules}

\subsubsection{Cosine Annealing}

\begin{equation}
\eta_t = \eta_{\min} + \frac{1}{2}(\eta_{\max} - \eta_{\min})\left(1 + \cos\left(\frac{t}{T}\pi\right)\right),
\label{eq:cosine_schedule}
\end{equation}

where $t$ is the current step and $T$ is the total number of steps.

\subsubsection{Warmup}

Linear warmup for the first $T_{\text{warmup}}$ steps:

\begin{equation}
\eta_t = \begin{cases}
\frac{t}{T_{\text{warmup}}} \eta_{\max}, & t \leq T_{\text{warmup}}, \\
\eta_t^{\text{(schedule)}}, & t > T_{\text{warmup}}.
\end{cases}
\label{eq:warmup}
\end{equation}

\chapter{Model Architecture}

This chapter describes the conditional U-Net architecture employed in our compact diffusion model, with emphasis on design choices for parameter efficiency and computational performance.

\section{U-Net Architecture Overview}

The U-Net~\citep{ronneberger2015u} is an encoder-decoder architecture with skip connections, originally designed for image segmentation but widely adopted for diffusion models~\citep{ho2020denoising}.

\begin{figure}[ht]
    \centering
    \includegraphics[width=0.9\linewidth]{U-net4.png}
    \caption{Complete U-Net architecture showing encoder, bottleneck, and decoder paths with skip connections.}
    \label{fig:unet_complete}
\end{figure}

\subsection{Encoder-Decoder Structure}

\subsubsection{Encoder}

The encoder progressively downsamples spatial dimensions while increasing channel capacity:

\begin{equation}
\mathbf{h}^{(\ell)} = \text{Down}^{(\ell)}(\mathbf{h}^{(\ell-1)}, \mathbf{e}_t),
\label{eq:encoder}
\end{equation}

where $\mathbf{h}^{(\ell)} \in \R^{C_\ell \times H_\ell \times W_\ell}$ with $H_\ell = H_{\ell-1}/2$ and $W_\ell = W_{\ell-1}/2$.

\subsubsection{Decoder}

The decoder upsamples features, incorporating skip connections from corresponding encoder layers:

\begin{equation}
\mathbf{h}^{(\ell)} = \text{Up}^{(\ell)}(\mathbf{h}^{(\ell+1)}, \mathbf{h}_{\text{skip}}^{(L-\ell)}, \mathbf{e}_t),
\label{eq:decoder}
\end{equation}

where $\mathbf{h}_{\text{skip}}^{(L-\ell)}$ is the corresponding encoder feature.

\section{Building Blocks}

\subsection{Double Convolution Block}

The fundamental building block is a sequence of two convolutions with normalization and activation:

\begin{align}
\mathbf{z}_1 &= \text{GELU}(\text{GroupNorm}(\text{Conv}(\mathbf{x}))), \label{eq:double_conv_1} \\
\mathbf{z}_2 &= \text{GELU}(\text{GroupNorm}(\text{Conv}(\mathbf{z}_1))). \label{eq:double_conv_2}
\end{align}

With residual connection:

\begin{equation}
\mathbf{y} = \mathbf{x} + \mathbf{z}_2,
\label{eq:double_conv_residual}
\end{equation}

where channel dimensions are matched via $1 \times 1$ convolution if necessary.

\begin{figure}[ht]
    \centering
    \includegraphics[width=0.4\linewidth]{Resnet-block.png}
    \caption{ResNet-style convolutional block with skip connection.}
    \label{fig:resnet_block}
\end{figure}

\subsection{Self-Attention Layer}

Self-attention enables modeling long-range dependencies:

\begin{equation}
\text{Attention}(\mathbf{Q}, \mathbf{K}, \mathbf{V}) = \text{softmax}\left(\frac{\mathbf{Q}\mathbf{K}^\top}{\sqrt{d_k}}\right)\mathbf{V},
\label{eq:attention}
\end{equation}

where $\mathbf{Q}, \mathbf{K}, \mathbf{V}$ are linear projections of the input features.

\begin{figure}[ht]
    \centering
    \includegraphics[width=0.3\linewidth]{self-attention-block.png}
    \caption{Self-attention block with layer normalization and feed-forward network.}
    \label{fig:attention_block}
\end{figure}

\subsubsection{Multi-Head Attention}

Split into $h$ heads:

\begin{equation}
\text{MultiHead}(\mathbf{X}) = \text{Concat}(\text{head}_1, \ldots, \text{head}_h)\mathbf{W}^O,
\label{eq:multihead}
\end{equation}

where $\text{head}_i = \text{Attention}(\mathbf{X}\mathbf{W}_i^Q, \mathbf{X}\mathbf{W}_i^K, \mathbf{X}\mathbf{W}_i^V)$.

\subsection{Downsampling and Upsampling}

\subsubsection{Downsampling}

Combine max pooling with strided convolution:

\begin{align}
\mathbf{h}_{\text{pool}} &= \text{MaxPool}(\mathbf{h}^{(\ell-1)}), \label{eq:maxpool} \\
\mathbf{h}^{(\ell)} &= \text{DoubleConv}(\mathbf{h}_{\text{pool}}, \mathbf{e}_t). \label{eq:down_conv}
\end{align}

\begin{figure}[ht]
    \centering
    \includegraphics[width=0.5\linewidth]{up-sample block.png}
    \caption{Upsampling block with bilinear interpolation and convolution.}
    \label{fig:upsample_block}
\end{figure}

\subsubsection{Upsampling}

Use bilinear interpolation followed by convolution:

\begin{align}
\mathbf{h}_{\text{up}} &= \text{Upsample}(\mathbf{h}^{(\ell+1)}, \text{scale}=2), \label{eq:upsample} \\
\mathbf{h}_{\text{concat}} &= \text{Concat}(\mathbf{h}_{\text{up}}, \mathbf{h}_{\text{skip}}^{(L-\ell)}), \label{eq:concat_skip} \\
\mathbf{h}^{(\ell)} &= \text{DoubleConv}(\mathbf{h}_{\text{concat}}, \mathbf{e}_t). \label{eq:up_conv}
\end{align}

\section{Conditioning Mechanisms}

\subsection{Timestep Embedding}

The diffusion timestep $t$ is embedded using sinusoidal functions:

\begin{equation}
\text{PE}(t, 2i) = \sin\left(\frac{t}{10000^{2i/d}}\right), \quad \text{PE}(t, 2i+1) = \cos\left(\frac{t}{10000^{2i/d}}\right),
\label{eq:timestep_embed}
\end{equation}

where $d$ is the embedding dimension.

These embeddings are passed through MLPs to match feature dimensions:

\begin{equation}
\mathbf{e}_t = \text{GELU}(\mathbf{W}_2 \text{GELU}(\mathbf{W}_1 \text{PE}(t))).
\label{eq:time_mlp}
\end{equation}

\subsection{Class Conditioning}

For class-conditional generation, class labels $y \in \{1, \ldots, C\}$ are embedded:

\begin{equation}
\mathbf{e}_y = \text{Embedding}(y),
\label{eq:class_embed}
\end{equation}

and added to the timestep embedding:

\begin{equation}
\mathbf{e} = \mathbf{e}_t + \mathbf{e}_y.
\label{eq:combined_embed}
\end{equation}

\subsection{Adaptive Group Normalization (AdaGN)}

Conditioning information modulates normalization statistics:

\begin{equation}
\text{AdaGN}(\mathbf{h}, \mathbf{e}) = \gamma(\mathbf{e}) \cdot \frac{\mathbf{h} - \mu}{\sigma} + \beta(\mathbf{e}),
\label{eq:adagn}
\end{equation}

where $\gamma$ and $\beta$ are predicted from $\mathbf{e}$ via linear layers.

\section{Compact Design Choices}

\subsection{Channel Compression Factor}

Introduce compression factor $c \in (0, 1]$:

\begin{equation}
C_\ell^{\text{compact}} = \left\lfloor c \cdot C_\ell^{\text{base}} \right\rfloor,
\label{eq:channel_compress}
\end{equation}

reducing parameter count by approximately $c^2$ (due to convolutional weight matrices).

\subsection{Attention Placement}

Attention is computationally expensive ($O(N^2)$ where $N = H \times W$). Apply attention only at low resolutions:

\begin{itemize}
\item Resolution $\geq 32 \times 32$: No attention
\item Resolution $16 \times 16$: Single attention layer
\item Resolution $8 \times 8$: Multiple attention layers
\end{itemize}

\subsection{Reduced Attention Heads}

Use fewer attention heads ($h = 4$ instead of $h = 8$) in compact models.

\subsection{Depthwise Separable Convolutions}

Replace standard convolutions with depthwise separable convolutions:

\begin{align}
\text{DepthwiseConv}(\mathbf{x}) &= \text{Conv}_{3 \times 3}^{\text{depthwise}}(\mathbf{x}), \label{eq:depthwise} \\
\text{PointwiseConv}(\mathbf{y}) &= \text{Conv}_{1 \times 1}(\mathbf{y}), \label{eq:pointwise}
\end{align}

reducing parameters from $C_{\text{in}} \cdot C_{\text{out}} \cdot k^2$ to $C_{\text{in}} \cdot k^2 + C_{\text{in}} \cdot C_{\text{out}}$.

\chapter{Symmetry-Aware Diffusion Models}

This chapter develops a comprehensive framework for incorporating physical symmetries and geometric invariances into diffusion model architectures. We draw from principles in group theory, geometric deep learning~\citep{bronstein2021geometric}, and physics-informed machine learning.

\section{Motivation and Mathematical Foundations}

\subsection{Symmetries in Physical Systems}

Physical systems often exhibit symmetries corresponding to transformations that leave the system's properties invariant. Examples include:

\begin{description}
\item[Translation Invariance] Physics laws are the same everywhere in space.
\item[Rotational Invariance] Isotropic systems have the same properties in all directions.
\item[Permutation Symmetry] Identical particles are indistinguishable.
\item[Gauge Symmetry] Physical observables are invariant under local phase transformations.
\end{description}

Incorporating these symmetries as inductive biases offers several advantages:

\begin{enumerate}
\item \textbf{Data Efficiency}: Fewer training examples required to learn symmetry-respecting functions.
\item \textbf{Generalization}: Better performance on out-of-distribution data respecting the same symmetries.
\item \textbf{Interpretability}: Generated samples automatically satisfy physical constraints.
\item \textbf{Parameter Efficiency}: Weight sharing induced by equivariance reduces model capacity needs.
\end{enumerate}

\subsection{Group Theory Fundamentals}

\begin{definition}[Group]
A group $(G, \cdot)$ is a set $G$ equipped with a binary operation $\cdot: G \times G \to G$ satisfying:
\begin{enumerate}
\item Closure: $\forall g_1, g_2 \in G$, $g_1 \cdot g_2 \in G$.
\item Associativity: $(g_1 \cdot g_2) \cdot g_3 = g_1 \cdot (g_2 \cdot g_3)$.
\item Identity: $\exists e \in G$ such that $e \cdot g = g \cdot e = g$ for all $g \in G$.
\item Inverse: $\forall g \in G$, $\exists g^{-1} \in G$ such that $g \cdot g^{-1} = g^{-1} \cdot g = e$.
\end{enumerate}
\end{definition}

\begin{definition}[Group Action]
A (left) action of group $G$ on set $X$ is a function $\rho: G \times X \to X$ satisfying:
\begin{enumerate}
\item $\rho(e, x) = x$ for all $x \in X$,
\item $\rho(g_1, \rho(g_2, x)) = \rho(g_1 \cdot g_2, x)$ for all $g_1, g_2 \in G, x \in X$.
\end{enumerate}
We often write $\rho(g, x) = g \cdot x$ or $T_g(x)$ where $T_g: X \to X$ is the transformation induced by $g$.
\end{definition}

\subsection{Equivariance and Invariance}

\begin{definition}[Equivariance]
A function $f: X \to Y$ is equivariant with respect to group actions $\rho_X$ on $X$ and $\rho_Y$ on $Y$ if:
\begin{equation}
f(\rho_X(g, x)) = \rho_Y(g, f(x)), \quad \forall g \in G, x \in X.
\label{eq:equivariance}
\end{equation}
Equivalently: $f \circ T_g^X = T_g^Y \circ f$.
\end{definition}

\begin{definition}[Invariance]
A function $f: X \to Y$ is invariant if $\rho_Y$ is the trivial action (identity):
\begin{equation}
f(\rho_X(g, x)) = f(x), \quad \forall g \in G, x \in X.
\label{eq:invariance}
\end{equation}
\end{definition}

\subsection{Relevant Symmetry Groups}

\subsubsection{Translation Group $\R^d$}

Group operation: vector addition. Action on functions:
\begin{equation}
(T_{\mathbf{v}} f)(\mathbf{x}) = f(\mathbf{x} - \mathbf{v}).
\label{eq:translation_action}
\end{equation}

Convolutional networks are naturally translation equivariant.

\subsubsection{Rotation Group $SO(d)$}

Special orthogonal group in $d$ dimensions. Action on points:
\begin{equation}
T_R(\mathbf{x}) = R\mathbf{x}, \quad R \in SO(d).
\label{eq:rotation_action}
\end{equation}

\subsubsection{Euclidean Group $E(d) = \R^d \rtimes SO(d)$}

Semidirect product of translations and rotations:
\begin{equation}
T_{(\mathbf{v}, R)}(\mathbf{x}) = R\mathbf{x} + \mathbf{v}.
\label{eq:euclidean_action}
\end{equation}

\subsubsection{Permutation Group $S_N$}

Group of permutations on $N$ elements. Action on sequences:
\begin{equation}
T_\pi(\mathbf{x}_1, \ldots, \mathbf{x}_N) = (\mathbf{x}_{\pi(1)}, \ldots, \mathbf{x}_{\pi(N)}).
\label{eq:permutation_action}
\end{equation}

Relevant for point clouds, molecular systems, particle physics.

\section{Equivariant Neural Networks}

\subsection{Group Convolution}

\begin{definition}[Group Convolution]
For group $G$ acting on feature space $X$, the group convolution is:
\begin{equation}
[f \star \psi](g) = \int_G f(h) \psi(h^{-1}g) \, dh,
\label{eq:group_conv}
\end{equation}
where $f: G \to \R^{C_{\text{in}}}$ is the input feature map and $\psi: G \to \R^{C_{\text{in}} \times C_{\text{out}}}$ is the filter.
\end{definition}

For discrete groups, the integral becomes a sum.

\subsection{Steerable Convolutions}

For continuous groups like $SO(2)$ or $SO(3)$, filters can be constructed from steerable bases~\citep{weiler2018learning}:

\begin{equation}
\psi(\mathbf{r}, R) = \sum_{\ell=0}^{L_{\max}} \sum_{m=-\ell}^\ell c_{\ell m} Y_\ell^m(R\mathbf{r}),
\label{eq:steerable_filter}
\end{equation}

where $Y_\ell^m$ are spherical harmonics and $c_{\ell m}$ are learnable coefficients.

\subsection{$E(3)$-Equivariant Graph Networks}

For molecular and particle data represented as graphs, Satorras et al.~\citep{satorras2021n} proposed $E(n)$-equivariant graph networks:

\begin{align}
\mathbf{m}_{ij} &= \phi_m(\mathbf{h}_i, \mathbf{h}_j, \norm{\mathbf{x}_i - \mathbf{x}_j}^2), \label{eq:egnn_message} \\
\mathbf{x}_i' &= \mathbf{x}_i + \sum_{j \neq i} (\mathbf{x}_i - \mathbf{x}_j) \phi_x(\mathbf{m}_{ij}), \label{eq:egnn_coord} \\
\mathbf{h}_i' &= \phi_h(\mathbf{h}_i, \sum_{j \neq i} \mathbf{m}_{ij}), \label{eq:egnn_feature}
\end{align}

where $\phi_m, \phi_x, \phi_h$ are MLPs, $\mathbf{x}_i \in \R^3$ are coordinates, and $\mathbf{h}_i \in \R^{C}$ are features.

Crucially, the coordinate update (Equation~\ref{eq:egnn_coord}) is equivariant:
\begin{equation}
T_R(\mathbf{x}_i') = R\mathbf{x}_i + \sum_{j \neq i} R(\mathbf{x}_i - \mathbf{x}_j) \phi_x(\mathbf{m}_{ij}) = R\mathbf{x}_i'.
\label{eq:egnn_equivariance_proof}
\end{equation}

\section{Symmetry-Aware Diffusion Models}

\subsection{Equivariant Forward Process}

For data $\mathbf{x}_0$ transforming under group $G$, the forward diffusion should respect symmetry:

\begin{equation}
q(T_g(\mathbf{x}_t) | T_g(\mathbf{x}_{t-1})) = q(\mathbf{x}_t | \mathbf{x}_{t-1}).
\label{eq:equivariant_forward}
\end{equation}

For Gaussian diffusion:
\begin{equation}
q(\mathbf{x}_t | \mathbf{x}_{t-1}) = \N(\mathbf{x}_t; \sqrt{1-\beta_t}\mathbf{x}_{t-1}, \beta_t\mathbf{I}),
\label{eq:gaussian_forward_symmetric}
\end{equation}

this is automatically satisfied for $G = \R^d$ (translations) and $G = SO(d)$ (rotations), since the Gaussian is isotropic.

\subsection{Equivariant Denoising Network}

The key requirement is that the denoising function $\boldsymbol{\epsilon}_\theta(\mathbf{x}_t, t)$ is equivariant:

\begin{equation}
\boldsymbol{\epsilon}_\theta(T_g(\mathbf{x}_t), t) = T_g(\boldsymbol{\epsilon}_\theta(\mathbf{x}_t, t)).
\label{eq:equivariant_denoiser}
\end{equation}

\subsubsection{Implementation via Equivariant Architectures}

Replace standard U-Net layers with group-equivariant analogs:

\begin{itemize}
\item \textbf{Translation Equivariance}: Standard convolutions.
\item \textbf{Rotation Equivariance}: Use steerable CNNs~\citep{weiler2018learning} or vector field networks~\citep{cohen2016group}.
\item \textbf{Permutation Equivariance}: Use graph networks or DeepSets~\citep{zaheer2017deep}.
\item \textbf{$E(3)$ Equivariance}: Use EGNN~\citep{satorras2021n} or Tensor Field Networks~\citep{thomas2018tensor}.
\end{itemize}

\subsection{Score Function Equivariance}

From the score-based perspective, equivariance of $\nabla_{\mathbf{x}} \log p(\mathbf{x})$ is natural:

\begin{proposition}[Score Equivariance]
If $p(T_g(\mathbf{x})) = p(\mathbf{x})$ for all $g \in G$ (density is invariant), then:
\begin{equation}
\nabla_{\mathbf{x}} \log p(T_g(\mathbf{x})) = J_g^\top \nabla_{T_g(\mathbf{x})} \log p(T_g(\mathbf{x})),
\label{eq:score_equivariance}
\end{equation}
where $J_g = \frac{\partial T_g}{\partial \mathbf{x}}$ is the Jacobian.
\end{proposition}

For linear actions ($T_g(\mathbf{x}) = g\mathbf{x}$ with $g$ a matrix), this simplifies to:
\begin{equation}
\nabla_{\mathbf{x}} \log p(g\mathbf{x}) = g^\top \nabla_{g\mathbf{x}} \log p(g\mathbf{x}).
\label{eq:score_linear_equivariance}
\end{equation}

\subsection{Gauge Equivariance}

For systems with gauge symmetry (e.g., molecular conformations invariant to global rotations), we can parameterize in a gauge-invariant coordinate system.

\subsubsection{Internal Coordinates}

For molecular systems, use bond lengths, angles, and dihedrals:
\begin{align}
d_{ij} &= \norm{\mathbf{x}_i - \mathbf{x}_j}, \label{eq:bond_length} \\
\theta_{ijk} &= \arccos\left(\frac{(\mathbf{x}_i - \mathbf{x}_j) \cdot (\mathbf{x}_k - \mathbf{x}_j)}{d_{ij} d_{jk}}\right), \label{eq:bond_angle} \\
\phi_{ijkl} &= \text{dihedral}(\mathbf{x}_i, \mathbf{x}_j, \mathbf{x}_k, \mathbf{x}_l). \label{eq:dihedral}
\end{align}

Diffusion operates in this internal coordinate space, ensuring $SE(3)$ invariance automatically.

\section{Incorporating Conservation Laws}

\subsection{Continuous Symmetries and Noether's Theorem}

Noether's theorem states that every continuous symmetry corresponds to a conserved quantity:

\begin{theorem}[Noether's Theorem]
If a Lagrangian $\mathcal{L}$ is invariant under a continuous transformation parameterized by $\epsilon$:
\begin{equation}
\mathbf{x}(t) \to \mathbf{x}(t) + \epsilon \delta\mathbf{x}(t),
\label{eq:infinitesimal_transform}
\end{equation}
then there exists a conserved charge:
\begin{equation}
Q = \frac{\partial \mathcal{L}}{\partial \dot{\mathbf{x}}} \cdot \delta\mathbf{x}.
\label{eq:conserved_charge}
\end{equation}
\end{theorem}

Examples:
\begin{itemize}
\item Translation symmetry $\Rightarrow$ momentum conservation.
\item Rotation symmetry $\Rightarrow$ angular momentum conservation.
\item Time translation symmetry $\Rightarrow$ energy conservation.
\end{itemize}

\subsection{Constrained Diffusion}

To enforce conservation laws, we can modify the diffusion process to operate on a constrained manifold~\citep{de2021riemannian}.

\subsubsection{Projection Method}

After each diffusion step, project onto the constraint manifold:

\begin{equation}
\mathbf{x}_t' = \mathop{\arg\min}_{\mathbf{x} \in \mathcal{M}} \norm{\mathbf{x} - \mathbf{x}_t}^2,
\label{eq:projection}
\end{equation}

where $\mathcal{M} = \{\mathbf{x} : Q(\mathbf{x}) = Q_0\}$ is the constraint manifold.

\subsubsection{Geodesic Diffusion}

For constraints defining a Riemannian manifold, diffusion can occur along geodesics:

\begin{equation}
d\mathbf{x} = -\nabla_{\mathcal{M}} \log p_t(\mathbf{x}) \, dt + \sqrt{2} \, d\mathbf{w}_{\mathcal{M}},
\label{eq:geodesic_diffusion}
\end{equation}

where $\nabla_{\mathcal{M}}$ is the Riemannian gradient and $d\mathbf{w}_{\mathcal{M}}$ is Brownian motion on $\mathcal{M}$.

\subsection{Energy-Based Models and Hamiltonian Dynamics}

For systems governed by Hamiltonian dynamics $H(\mathbf{q}, \mathbf{p})$, we can define diffusion in phase space:

\begin{align}
d\mathbf{q} &= \frac{\partial H}{\partial \mathbf{p}} \, dt + \sigma_q \, d\mathbf{w}_q, \label{eq:hamiltonian_q} \\
d\mathbf{p} &= -\frac{\partial H}{\partial \mathbf{q}} \, dt + \sigma_p \, d\mathbf{w}_p. \label{eq:hamiltonian_p}
\end{align}

The denoising network learns:
\begin{equation}
\boldsymbol{\epsilon}_\theta(\mathbf{q}, \mathbf{p}, t) = \left(\frac{\partial H_\theta}{\partial \mathbf{p}}, -\frac{\partial H_\theta}{\partial \mathbf{q}}\right),
\label{eq:hamiltonian_score}
\end{equation}

ensuring generated trajectories respect energy conservation (up to noise).

\section{Symmetry-Induced Model Compression}

\subsection{Parameter Sharing via Equivariance}

Equivariant networks naturally share parameters across group orbits, reducing effective capacity.

\begin{proposition}[Parameter Count Reduction]
For a group $G$ acting on input/output spaces, a $G$-equivariant linear layer has parameter count:
\begin{equation}
|\Theta_{\text{equiv}}| = \frac{|\Theta_{\text{standard}}|}{|G|},
\label{eq:parameter_reduction}
\end{equation}
where $|G|$ is the group order (for finite groups).
\end{proposition}

For continuous groups, the reduction factor depends on the representation dimensionality.

\subsection{Factorized Generation}

For systems with independent symmetry sectors, generation can be factorized:

\begin{equation}
p(\mathbf{x}) = \prod_{i=1}^K p(\mathbf{x}^{(i)}),
\label{eq:factorized_generation}
\end{equation}

where $\mathbf{x}^{(i)}$ are disjoint components transformed by subgroups $G_i \subseteq G$.

Train separate compact models for each sector:
\begin{equation}
\boldsymbol{\epsilon}_\theta(\mathbf{x}_t, t) = \bigoplus_{i=1}^K \boldsymbol{\epsilon}_{\theta_i}(\mathbf{x}_t^{(i)}, t).
\label{eq:factorized_denoiser}
\end{equation}

Total parameter count: $\sum_i |\theta_i| < |\theta_{\text{monolithic}}|$.

\section{Examples and Applications}

\subsection{Image Generation with Rotation Equivariance}

For natural images exhibiting approximate rotational symmetry (e.g., flowers, wheels):

Use $SO(2)$-equivariant convolutions~\citep{weiler2018learning}. Filter bank:
\begin{equation}
\psi^{(k)}(r, \theta) = \sum_{\ell=0}^{L_{\max}} c_{\ell}^{(k)} R_\ell(r) e^{i\ell\theta},
\label{eq:so2_filter}
\end{equation}

where $R_\ell(r)$ are radial functions and $c_{\ell}^{(k)}$ are learned coefficients.

\subsection{Molecular Conformation Generation}

For molecule generation respecting $SE(3)$ symmetry~\citep{hoogeboom2022equivariant}:

\begin{enumerate}
\item Represent molecule as graph: nodes = atoms, edges = bonds.
\item Use EGNN layers (Equations~\ref{eq:egnn_message}--\ref{eq:egnn_feature}) for $\boldsymbol{\epsilon}_\theta$.
\item Diffusion operates on 3D coordinates $\{\mathbf{x}_i\}$ and node features $\{\mathbf{h}_i\}$.
\item Conservation: enforce center of mass at origin, optionally constrain bond lengths.
\end{enumerate}

\subsection{Point Cloud Generation}

For 3D point clouds (e.g., LiDAR, CAD):

Use permutation-equivariant architectures like PointNet++~\citep{qi2017pointnet++} or DGCNN~\citep{wang2019dynamic} as the backbone for $\boldsymbol{\epsilon}_\theta$.

Diffusion on point clouds:
\begin{equation}
\mathbf{P}_t = \{\mathbf{x}_i^{(t)}\}_{i=1}^N, \quad \mathbf{x}_i^{(t)} = \sqrt{\bar{\alpha}_t}\mathbf{x}_i^{(0)} + \sqrt{1-\bar{\alpha}_t}\boldsymbol{\epsilon}_i.
\label{eq:point_cloud_diffusion}
\end{equation}

Denoising network must be permutation equivariant:
\begin{equation}
\boldsymbol{\epsilon}_\theta(T_\pi(\mathbf{P}_t), t) = T_\pi(\boldsymbol{\epsilon}_\theta(\mathbf{P}_t, t)).
\label{eq:permutation_equivariant_denoiser}
\end{equation}

\subsection{Physics Simulation}

For particle systems governed by pairwise potentials $V(\mathbf{x}) = \sum_{i<j} v(\norm{\mathbf{x}_i - \mathbf{x}_j})$:

Score function:
\begin{equation}
\nabla_{\mathbf{x}_i} \log p(\mathbf{x}) = -\nabla_{\mathbf{x}_i} V(\mathbf{x}) = -\sum_{j \neq i} v'(\norm{\mathbf{x}_i - \mathbf{x}_j}) \frac{\mathbf{x}_i - \mathbf{x}_j}{\norm{\mathbf{x}_i - \mathbf{x}_j}}.
\label{eq:pairwise_score}
\end{equation}

This is naturally $E(3)$-equivariant. Neural network learns the potential:
\begin{equation}
v_\theta(r) \approx v(r).
\label{eq:learned_potential}
\end{equation}

\section{Combining Symmetry with Other Compression Techniques}

\subsection{Symmetry-Aware Pruning}

Pruning must respect equivariance constraints. For group-equivariant layers, prune entire irreducible representations (irreps) rather than individual channels:

\begin{equation}
\text{Prune irrep } \rho_\ell \text{ if } \text{importance}(\rho_\ell) < \delta,
\label{eq:irrep_pruning}
\end{equation}

where importance aggregates across all channels in the irrep.

\subsection{Symmetry-Aware Quantization}

Quantization scales can be shared across group orbits:

\begin{equation}
\Delta_g = \Delta \quad \text{for all } g \in G,
\label{eq:shared_quantization_scale}
\end{equation}

reducing the number of quantization parameters.

\subsection{Distillation with Symmetry Constraints}

Student and teacher must both respect the same symmetries. Distillation loss can include symmetry-specific terms:

\begin{equation}
\loss_{\text{sym}} = \E_{g \sim G}\left[\norm{\boldsymbol{\epsilon}_S(T_g(\mathbf{x}_t), t) - T_g(\boldsymbol{\epsilon}_S(\mathbf{x}_t, t))}^2\right],
\label{eq:symmetry_distillation_loss}
\end{equation}

explicitly penalizing symmetry violations.

\chapter{Experimental Results and Analysis}

This chapter presents our experimental methodology, quantitative results, and critical analysis of the compact diffusion models developed in this work.

\section{Experimental Setup}

\subsection{Dataset}

We use the CIFAR-10 dataset~\citep{krizhevsky2009learning}, consisting of 60,000 color images ($32 \times 32$ pixels) across 10 classes. Due to computational constraints:

\begin{itemize}
\item \textbf{Training Set}: 10,000 randomly sampled images.
\item \textbf{Validation Set}: 1,000 images.
\item \textbf{Test Set}: 1,000 images.
\end{itemize}

Images are resized to $64 \times 64$ pixels and normalized to $[-1, 1]$.

\subsection{Model Configuration}

\subsubsection{Teacher Model}

\begin{itemize}
\item \textbf{Architecture}: U-Net with base channels $C = 128$, compression factor $c = 1$.
\item \textbf{Timesteps}: $T = 1000$, linear variance schedule $\beta_1 = 10^{-4}$ to $\beta_T = 0.02$.
\item \textbf{Self-Attention}: At resolutions $16 \times 16$ and $8 \times 8$, 4 heads.
\end{itemize}

\subsubsection{Student Model}

\begin{itemize}
\item \textbf{Architecture}: U-Net with base channels $C = 64$, compression factor $c = 0.5$.
\item \textbf{Distillation}: Combined loss with $\alpha = 0.5$, $\beta = 0.5$ (Equation~\ref{eq:combined_distill_loss}).
\item \textbf{Feature Matching}: 3 intermediate layers, $\gamma_\ell = 0.1$.
\end{itemize}

\subsubsection{Pruned Student Model}

\begin{itemize}
\item \textbf{Pruning Method}: Structured channel pruning with L1 norm importance.
\item \textbf{Pruning Ratio}: Global 30\% sparsity.
\item \textbf{Fine-Tuning}: 50 epochs post-pruning.
\end{itemize}

\subsection{Training Details}

\begin{itemize}
\item \textbf{Optimizer}: AdamW with $\beta_1 = 0.9$, $\beta_2 = 0.999$, weight decay $\lambda = 10^{-4}$.
\item \textbf{Learning Rate}: Fixed $\eta = 10^{-5}$ (suboptimal; see discussion).
\item \textbf{Batch Size}: 64 (with gradient accumulation to simulate batch size 128).
\item \textbf{Epochs}: Approximately 200 (limited by compute).
\item \textbf{Mixed Precision}: FP16 with dynamic loss scaling.
\item \textbf{EMA}: Decay $\beta = 0.995$, applied after 1000 warmup steps.
\end{itemize}

\subsection{Hardware}

Single NVIDIA V100 GPU (16GB VRAM), training time approximately 48 hours for teacher, 36 hours for student.

\section{Quantitative Results}

\subsection{Model Size and Inference Time}

\begin{table}[H]
\centering
\begin{tabular}{lccc}
\toprule
\textbf{Model} & \textbf{Parameters (M)} & \textbf{Size (MB)} & \textbf{Inference Time (s/batch)} \\
\midrule
Teacher & 23.5 & 89.02 & 18.54 \\
Student (Distilled) & 6.0 & 22.60 & 15.72 \\
Student (Pruned 30\%) & 4.6 & 17.42 & 15.40 \\
\bottomrule
\end{tabular}
\caption{Model size and inference time comparison. Batch size = 8, $T = 1000$ sampling steps.}
\label{tab:model_sizes}
\end{table}

\textbf{Key Observations}:
\begin{itemize}
\item Student achieves 74.6\% size reduction compared to teacher.
\item Pruning provides additional 23\% reduction (15\% relative to student).
\item Inference speedup modest (15\%) due to memory-bound operations; more significant speedup expected on mobile/edge devices.
\end{itemize}

\subsection{Generation Quality: FID Scores}

\begin{table}[H]
\centering
\begin{tabular}{lccc}
\toprule
\textbf{Model} & \textbf{FID} $\downarrow$ & \textbf{Inception Score} $\uparrow$ & \textbf{Training Epochs} \\
\midrule
Teacher & 163.2 & 3.41 & 200 \\
Student (Distilled) & 184.7 & 3.12 & 200 \\
Student (Pruned 30\%) & 213.5 & 2.85 & 200 + 50 fine-tune \\
\bottomrule
\end{tabular}
\caption{Generation quality metrics. Lower FID is better; higher Inception Score is better. For reference, fully-trained DDPM on full CIFAR-10 achieves FID $\approx$ 3--5.}
\label{tab:fid_scores}
\end{table}

\textbf{Key Observations}:
\begin{itemize}
\item All models exhibit elevated FID scores, indicating significant deviation from real data distribution.
\item Student suffers 13\% FID degradation relative to teacher.
\item Aggressive pruning causes additional 16\% FID degradation.
\end{itemize}

\subsection{Training Curves}

\begin{figure}[H]
\centering
\includegraphics[width=0.6\linewidth]{teacher-loss.png}
\caption{MSE training loss for teacher model over 200 epochs. Loss decreases steadily but has not fully converged.}
\label{fig:teacher_loss}
\end{figure}

\begin{figure}[H]
\centering
\includegraphics[width=0.6\linewidth]{student-model-loss.png}
\caption{MSE training loss for student model with knowledge distillation. Combined loss shown; standard diffusion loss and distillation loss behave similarly.}
\label{fig:student_loss}
\end{figure}

\subsection{Generated Samples}

\begin{figure}[H]
\centering
\includegraphics[width=0.6\linewidth]{teacher-model-images.png}
\caption{Sample images generated by the teacher model after 200 epochs. Recognizable class features but lacking fine details and exhibiting noise artifacts.}
\label{fig:teacher_images}
\end{figure}

\begin{figure}[H]
\centering
\includegraphics[width=0.8\linewidth]{trained-student-model-with_MPT.png}
\caption{Sample images generated by the distilled student model. Quality comparable to teacher but with slightly reduced sharpness.}
\label{fig:student_images}
\end{figure}

\begin{figure}[H]
\centering
\includegraphics[width=0.8\linewidth]{pruned-model-images.png}
\caption{Sample images generated by the pruned student model (30\% sparsity). Visible quality degradation with increased noise and loss of fine details.}
\label{fig:pruned_images}
\end{figure}

\section{Critical Analysis}

\subsection{Factors Contributing to High FID Scores}

\subsubsection{Insufficient Training}

Our models were trained for only 200 epochs on 10,000 images. State-of-the-art diffusion models typically require:
\begin{itemize}
\item \textbf{Epochs}: 500--1000+ epochs
\item \textbf{Data}: Full datasets (50,000--millions of images)
\item \textbf{Compute}: Multiple high-end GPUs for days/weeks
\end{itemize}

The training curves (Figures~\ref{fig:teacher_loss},~\ref{fig:student_loss}) show continued decrease, indicating lack of convergence.

\subsubsection{Fixed Learning Rate}

Using a fixed learning rate $\eta = 10^{-5}$ throughout training is suboptimal. Adaptive schedules (cosine annealing, warm restarts) enable:
\begin{itemize}
\item Initial rapid learning with higher rates
\item Fine-grained optimization with decreased rates near convergence
\item Escape from local minima via periodic rate increases
\end{itemize}

\textbf{Estimated Impact}: Could reduce FID by 30--50 points.

\subsubsection{Limited Dataset}

Training on 10,000 images (16.7\% of CIFAR-10) severely limits:
\begin{itemize}
\item Diversity of learned features
\item Generalization to unseen modes
\item Coverage of intra-class variations
\end{itemize}

Data augmentation (random crops, flips, color jitter) was minimal.

\textbf{Estimated Impact}: Could reduce FID by 40--60 points.

\subsubsection{Aggressive Compression}

The student model has 74\% fewer parameters than the teacher. While knowledge distillation mitigates capacity reduction, some performance loss is inevitable. The pruned model (30\% sparsity) removes critical connections, as evidenced by the FID degradation from 184.7 to 213.5.

\subsection{Comparison to Baselines}

\begin{table}[H]
\centering
\begin{tabular}{lccc}
\toprule
\textbf{Model} & \textbf{Dataset} & \textbf{Params (M)} & \textbf{FID} \\
\midrule
DDPM~\citep{ho2020denoising} & CIFAR-10 (full) & 35.7 & 3.17 \\
Improved DDPM~\citep{nichol2021improved} & CIFAR-10 (full) & 35.7 & 2.90 \\
Our Teacher & CIFAR-10 (10k subset) & 23.5 & 163.2 \\
Our Student & CIFAR-10 (10k subset) & 6.0 & 184.7 \\
\bottomrule
\end{tabular}
\caption{Comparison to published baselines. Note: direct comparison is unfair due to different training regimes.}
\label{tab:baselines}
\end{table}

The large FID gap (160+ points) is primarily attributable to training regime differences rather than architectural deficiencies.

\subsection{Distillation Effectiveness}

Despite the absolute FID scores being high, the \textit{relative} performance is informative:

\begin{itemize}
\item Student achieves 74\% size reduction with only 13\% FID degradation (184.7 vs 163.2).
\item This corresponds to a compression-quality trade-off of $\Delta\text{FID} / \Delta\text{size} \approx 0.24$ FID points per MB saved.
\end{itemize}

\textbf{Interpretation}: Knowledge distillation successfully transfers learned representations despite capacity constraints.

\subsection{Pruning Trade-offs}

Pruning 30\% of channels yields:
\begin{itemize}
\item \textbf{Benefit}: 23\% size reduction, minor inference speedup
\item \textbf{Cost}: 16\% FID degradation (213.5 vs 184.7)
\end{itemize}

\textbf{Analysis}: The FID increase indicates that the pruned channels contained non-redundant information. Less aggressive pruning (e.g., 15--20\%) or iterative pruning with longer fine-tuning would likely improve the trade-off.

\section{Ablation Studies}

\subsection{Impact of Distillation Loss Components}

\begin{table}[H]
\centering
\begin{tabular}{lccc}
\toprule
\textbf{Configuration} & \textbf{FID} & \textbf{$\alpha$} & \textbf{$\beta$} \\
\midrule
Standard diffusion loss only & 198.3 & 1.0 & 0.0 \\
Output distillation only & 192.1 & 0.0 & 1.0 \\
Combined ($\alpha = \beta = 0.5$) & 184.7 & 0.5 & 0.5 \\
Combined + feature matching & 182.4 & 0.5 & 0.5 \\
\bottomrule
\end{tabular}
\caption{Ablation study on distillation loss components. Feature matching provides small but consistent improvement.}
\label{tab:ablation_distillation}
\end{table}

\subsection{Impact of EMA}

\begin{table}[H]
\centering
\begin{tabular}{lcc}
\toprule
\textbf{Configuration} & \textbf{FID (raw weights)} & \textbf{FID (EMA weights)} \\
\midrule
Teacher & 171.5 & 163.2 \\
Student & 192.8 & 184.7 \\
\bottomrule
\end{tabular}
\caption{Effect of using EMA weights for evaluation. EMA consistently improves FID by 8--10 points.}
\label{tab:ablation_ema}
\end{table}

\section{Limitations}

\subsection{Computational Constraints}

\begin{itemize}
\item \textbf{Training Time}: Limited to 48 hours per model due to hardware availability.
\item \textbf{Hyperparameter Search}: No systematic search; parameters chosen based on prior work and manual tuning.
\item \textbf{Ensemble Models}: Could not train multiple runs for statistical significance.
\end{itemize}

\subsection{Architectural Choices}

\begin{itemize}
\item \textbf{Attention Layers}: Limited to low resolutions; full-resolution attention might improve quality.
\item \textbf{Depth}: Relatively shallow U-Net; deeper architectures could learn richer representations.
\item \textbf{Advanced Modules}: Did not explore recent innovations like cross-attention, adaptive layer norm, or flash attention.
\end{itemize}

\subsection{Evaluation Metrics}

\begin{itemize}
\item \textbf{FID}: Sensitive to sample size and feature extractor (Inception-v3); not always aligned with perceptual quality.
\item \textbf{Human Evaluation}: No user studies conducted; perceptual metrics (LPIPS) could provide complementary insights.
\item \textbf{Diversity}: Did not explicitly measure mode coverage or diversity (e.g., via precision/recall metrics).
\end{itemize}

\chapter{Conclusion and Future Directions}

\section{Summary of Contributions}

This work has presented a comprehensive theoretical and empirical study of compact diffusion models, with the following key contributions:

\begin{enumerate}
\item \textbf{Rigorous Mathematical Foundations}: Complete derivations of diffusion model training objectives, including variational lower bounds, score-based formulations, and SDE perspectives.

\item \textbf{Compression Techniques}: Detailed analysis of knowledge distillation, structured pruning, and quantization for diffusion models, with novel formulations for timestep-aware importance metrics and progressive distillation.

\item \textbf{Symmetry-Aware Framework}: Introduction of a comprehensive framework for incorporating physical symmetries into diffusion architectures, connecting geometric deep learning to generative modeling.

\item \textbf{Empirical Validation}: Demonstration of 74\% model size reduction with controlled performance degradation through knowledge distillation and pruning.
\end{enumerate}

\section{Key Insights}

\begin{itemize}
\item \textbf{Distillation Efficacy}: Knowledge distillation successfully transfers learned representations from teacher to student, maintaining functional performance (as measured by training loss) despite significant capacity reduction.

\item \textbf{Training Regime Criticality}: The quality of diffusion models is highly sensitive to training duration, learning rate schedules, and dataset size. Our elevated FID scores primarily reflect insufficient training rather than architectural deficiencies.

\item \textbf{Pruning Trade-offs}: Structured channel pruning provides marginal inference speedup but can significantly degrade generation quality if applied too aggressively. Iterative pruning with extensive fine-tuning is essential.

\item \textbf{Symmetry Benefits}: Incorporating symmetries as inductive biases offers a promising avenue for improving both parameter efficiency and data efficiency, though implementation requires careful architectural design.
\end{itemize}

\section{Future Research Directions}

\subsection{Improved Training Protocols}

\begin{enumerate}
\item \textbf{Extended Training}: Train on full CIFAR-10 (or larger datasets like ImageNet) for 500--1000 epochs.
\item \textbf{Adaptive Learning Rates}: Implement cosine annealing with warm restarts; explore learning rate warmup.
\item \textbf{Data Augmentation}: Extensive augmentation strategies (AutoAugment, RandAugment) to improve generalization with limited data.
\item \textbf{Large Batch Training}: Utilize gradient accumulation and distributed training for effective batch sizes of 256--512.
\end{enumerate}

\subsection{Advanced Compression}

\begin{enumerate}
\item \textbf{Quantization Implementation}: Complete quantization-aware training pipeline with per-channel quantization and dynamic bit-width allocation.
\item \textbf{Progressive Distillation}: Implement full progressive distillation~\citep{salimans2022progressive} to reduce sampling steps from 1000 to 4--8.
\item \textbf{Consistency Models}: Explore consistency distillation~\citep{song2023consistency} for single-step generation.
\item \textbf{Hybrid Methods}: Combine multiple compression techniques (e.g., quantization + pruning + distillation) with careful coordination.
\end{enumerate}

\subsection{Latent Diffusion Models}

Transition to latent diffusion~\citep{rombach2022high} for significantly improved efficiency:

\begin{itemize}
\item Train variational autoencoder (VAE) to compress images to low-dimensional latent space.
\item Perform diffusion in latent space (e.g., $4 \times 4 \times C$ instead of $64 \times 64 \times 3$).
\item Expected benefits: 8--16$\times$ speedup, improved sample quality, reduced memory footprint.
\end{itemize}

\subsection{Symmetry-Aware Architectures}

\begin{enumerate}
\item \textbf{$SO(2)$ Equivariant U-Nets}: Implement steerable CNNs for rotation-equivariant image generation.
\item \textbf{Point Cloud Diffusion}: Develop permutation-equivariant diffusion models for 3D point clouds using PointNet++ or DGCNN backbones.
\item \textbf{Molecular Generation}: Apply $E(3)$-equivariant graph networks for drug discovery and materials science applications.
\item \textbf{Physics Simulation}: Incorporate conservation laws (energy, momentum, angular momentum) via constrained diffusion on Riemannian manifolds.
\end{enumerate}

\subsection{Theoretical Analysis}

\begin{enumerate}
\item \textbf{Convergence Guarantees}: Formal analysis of convergence rates for distilled and pruned models.
\item \textbf{Capacity Bounds}: Determine minimum model capacity required to approximate teacher performance within $\epsilon$ error.
\item \textbf{Symmetry Breaking}: Study conditions under which learned representations spontaneously break or preserve symmetries.
\item \textbf{Score Approximation}: Error bounds for score function approximation in equivariant networks.
\end{enumerate}

\subsection{Application Domains}

\begin{enumerate}
\item \textbf{Medical Imaging}: Compact diffusion models for MRI/CT image synthesis with privacy-preserving constraints.
\item \textbf{Scientific Computing}: Accelerated sampling for molecular dynamics, weather forecasting, and computational fluid dynamics.
\item \textbf{Edge Deployment}: Mobile/IoT implementation with hardware-aware neural architecture search.
\item \textbf{Controllable Generation}: Integrate symmetry-aware diffusion with conditional generation (text-to-image, image editing).
\end{enumerate}

\section{Broader Impact}

Developing compact generative models has significant implications:

\subsection{Positive Impacts}

\begin{itemize}
\item \textbf{Democratization}: Enabling generative AI on consumer hardware, reducing barriers to entry.
\item \textbf{Energy Efficiency}: Reduced computational requirements lead to lower carbon footprint.
\item \textbf{Privacy}: On-device generation avoids data transmission to centralized servers.
\item \textbf{Scientific Discovery}: Accelerating simulation-based inference in physics, chemistry, and biology.
\end{itemize}

\subsection{Potential Concerns}

\begin{itemize}
\item \textbf{Misuse}: Compact models could facilitate creation of deepfakes or disinformation at scale.
\item \textbf{Bias Amplification}: Compressed models may inherit or amplify biases from training data.
\item \textbf{Accessibility}: Making generative models more accessible could exacerbate existing challenges around content authenticity.
\end{itemize}

Responsible development requires careful consideration of these trade-offs and implementation of appropriate safeguards.

\section{Concluding Remarks}

This work demonstrates that substantial compression of diffusion models is achievable through strategic application of knowledge distillation, pruning, and architectural innovations. While our experimental results are limited by computational constraints, the theoretical frameworks and methodologies developed here provide a foundation for future research.

The integration of symmetry-aware architectures represents a particularly promising direction, offering a principled approach to improving both parameter efficiency and physical consistency. As diffusion models continue to advance toward practical deployment, combining these compression techniques with domain-specific inductive biases will be essential for realizing their full potential across diverse application domains.

The path toward truly compact, efficient, and reliable generative models remains challenging, but the techniques and insights presented in this work contribute meaningful steps toward that goal.

\nocite{*}
\bibliographystyle{apalike}
\bibliography{references}

\appendix

\chapter{Additional Experimental Details}

\section{Hyperparameter Settings}

\begin{table}[H]
\centering
\begin{tabular}{lc}
\toprule
\textbf{Hyperparameter} & \textbf{Value} \\
\midrule
Batch size & 64 \\
Gradient accumulation steps & 2 \\
Learning rate & $10^{-5}$ \\
Optimizer & AdamW \\
$\beta_1$ & 0.9 \\
$\beta_2$ & 0.999 \\
Weight decay & $10^{-4}$ \\
EMA decay & 0.995 \\
Diffusion timesteps $T$ & 1000 \\
$\beta_{\min}$ & $10^{-4}$ \\
$\beta_{\max}$ & 0.02 \\
Mixed precision & FP16 \\
Loss scale & Dynamic \\
\bottomrule
\end{tabular}
\caption{Complete hyperparameter settings for all experiments.}
\end{table}

\section{Computational Requirements}

\begin{table}[H]
\centering
\begin{tabular}{lccc}
\toprule
\textbf{Model} & \textbf{Training Time (hrs)} & \textbf{GPU Memory (GB)} & \textbf{FLOPs/sample} \\
\midrule
Teacher & 48 & 14.2 & $2.1 \times 10^{12}$ \\
Student & 36 & 8.7 & $5.4 \times 10^{11}$ \\
Pruned Student & 42 (inc. fine-tune) & 7.1 & $4.2 \times 10^{11}$ \\
\bottomrule
\end{tabular}
\caption{Computational requirements for model training and inference.}
\end{table}

\end{document}
